\documentclass[11pt]{article}
\usepackage[a4paper,margin=21mm]{geometry}
\usepackage[no-math]{fontspec}
\setmainfont{Latin Modern Roman}
\newfontfamily\CJKfont{Noto Serif CJK SC}
\XeTeXlinebreaklocale "zh"
\XeTeXlinebreakskip=0pt plus 1pt
\usepackage{amsmath,amssymb,amsthm,amscd,graphicx,color}
\usepackage{xurl}
\usepackage[unicode,hidelinks]{hyperref}
\allowdisplaybreaks
\setlength\emergencystretch{3em}
\numberwithin{equation}{section}

\numberwithin{equation}{section}

\newtheorem{Theorem}{Theorem}[section]
\newtheorem*{Theorem*}{Theorem}
\newtheorem{Corollary}[Theorem]{Corollary}
\newtheorem{Lemma}[Theorem]{Lemma}
\newtheorem{Proposition}[Theorem]{Proposition}
{ \theoremstyle{definition}
	\newtheorem{Definition}[Theorem]{Definition}
	\newtheorem{Note}[Theorem]{Note}
	\newtheorem{Example}[Theorem]{Example}
	\newtheorem{Remark}[Theorem]{Remark} }


\usepackage{tikz}
\usetikzlibrary{matrix,arrows,patterns,cd}
\usepackage{leftidx}

\usepackage[mathscr]{eucal}

\newcommand{\Lbb}{{\boldsymbol{\EuScript{L}}}}
\newcommand{\Mbb}{{\boldsymbol{\EuScript{M}}}}
\newcommand{\Nbb}{{\boldsymbol{\EuScript{N}}}}
\newcommand{\Rbb}{{\boldsymbol{\EuScript{R}}}}
\newcommand{\Ubb}{{\boldsymbol{\EuScript{U}}}}

\newcommand{\II}{{\boldsymbol{1}}}

\newcommand{\CC}{{\mathbb C}}
\newcommand{\Kbb}{{\mathbb K}}
\newcommand{\RR}{{\mathbb R}}
\newcommand{\MM}{{\mathbb M}}
\newcommand{\NN}{{\mathbb N}}
%\newcommand{\QQ}{{\mathbb Q}}
\newcommand{\TT}{{\mathbb T}}
\newcommand{\ZZ}{{\mathbb Z}}


% Various classes of smooth compactly supported functions

\newcommand{\CoinfM}{C_0^\infty(M)}
\newcommand{\CoinfN}{C_0^\infty(N)}
\newcommand{\Coinfd}{C_0^\infty(\RR^d\backslash\{ 0\})}
\newcommand{\Coinf}[1]{C_0^\infty(\RR^{#1}\backslash\{ 0\})}
\newcommand{\CoinX}[1]{C_0^\infty\big({#1}\big)}
\newcommand{\GoinX}[1]{\Gamma_0^\infty({#1})}
\newcommand{\Coin}{C_0^\infty(0,\infty)}

\newcommand{\fc}{\mathrm{fc}}
\newcommand{\pc}{\mathrm{pc}}
\newcommand{\sfc}{\mathrm{sfc}}
\newcommand{\spc}{\mathrm{spc}}


% Sanserif
\newcommand{\sN}{{\mathsf N}}
\newcommand{\Rs}{{\mathsf R}}
\newcommand{\Ts}{{\mathsf T}}
\newcommand{\Ss}{{\mathsf S}}
\newcommand{\Sshat}{\widehat{\mathsf S}}
\newcommand{\fs}{{\mathsf f}}
\newcommand{\Js}{{\mathsf J}}
\newcommand{\vs}{{\mathsf v}}


% Calligraphic characters (N.B. to use rsfs use \mathscr)

\newcommand{\BB}{{\mathscr B}}
\newcommand{\DD}{{\mathscr D}}
\newcommand{\EE}{{\mathscr E}}
\newcommand{\HH}{{\mathscr H}}
\newcommand{\GG}{{\mathscr G}}
\newcommand{\KK}{{\mathscr K}}
\newcommand{\FF}{{\mathscr F}}
\newcommand{\LL}{{\mathcal L}}

\newcommand{\Ac}{{\mathcal A}}
\newcommand{\Bc}{{\mathcal B}}
\newcommand{\Cc}{{\mathcal C}}
\newcommand{\Dc}{{\mathcal D}}
\newcommand{\Fc}{{\mathcal F}}
\newcommand{\Gc}{{\mathcal G}}
\newcommand{\Hc}{{\mathcal H}}
\newcommand{\Jc}{{\mathcal J}}
\newcommand{\cN}{{\mathcal N}}
\newcommand{\cP}{{\mathcal P}}
\newcommand{\cR}{{\mathcal R}}
\newcommand{\cT}{{\mathcal T}}
\newcommand{\Uc}{{\mathcal U}}
\newcommand{\Xc}{{\mathcal X}}
\newcommand{\Zc}{{\mathcal Z}}

\newcommand{\Ai}{{\mathcal A}^\circ}
\newcommand{\cV}{{\mathcal V}}
\newcommand{\cI}{{\mathcal I}}
\newcommand{\Ic}{{\mathcal I}}
\newcommand{\cK}{{\mathcal K}}

\newcommand{\OO}{{\mathscr O}}

\newcommand{\QQ}{{\mathcal Q}}
\newcommand{\Aa}{{\cal A}}
\newcommand{\Ba}{{\cal B}}
\newcommand{\Fa}{{\cal F}}
%\newcommand{\FfL}{{\cal F}}
%\newcommand{\Ss}{{\cal S}}
\newcommand{\Uu}{{\cal U}}
\newcommand{\Xx}{{\cal X}}
\newcommand{\Zz}{{\cal Z}}
\newcommand{\Ft}{{\widetilde{\FfL}}}


\newcommand{\sW}{{\sf W}}


% Boldface characters

\newcommand{\kb}{{\boldsymbol{k}}}
\newcommand{\lb}{{\boldsymbol{\ell}}}
\newcommand{\fb}{{\boldsymbol{f}}}
\newcommand{\gb}{{\boldsymbol{g}}}
\newcommand{\hb}{{\boldsymbol{h}}}
\newcommand{\nb}{{\boldsymbol{n}}}
\newcommand{\pb}{{\vec{p}}}
\newcommand{\xb}{{\boldsymbol{x}}}
\newcommand{\yb}{{\boldsymbol{y}}}
\newcommand{\xbo}{{\boldsymbol{x_0}}}
\newcommand{\etb}{{\boldsymbol{\eta}}}
\newcommand{\Tb}{{\boldsymbol{T}}}

% Miscellaneous
\newcommand{\ab}{\mathrm{ab}}
\newcommand{\Afrak}{{\ygoth A}}
\newcommand{\Of}{{\mathfrak O}}
\newcommand{\ff}{{\mathfrak F}}
\newcommand{\Ip}{{\mathfrak I}} %%% integer part

\newcommand{\Afr}{{\mathfrak A}}
\newcommand{\Mfr}{{\mathfrak M}}
\newcommand{\Nfr}{{\mathfrak N}}

\newcommand{\ogth}{{\mathfrak o}}
\newcommand{\tgth}{{\mathfrak t}}
\newcommand{\wgth}{{\mathfrak w}}

\newcommand{\etap}{\eta'}
%\newcommand{\Dal}{\fbox{\phantom{${\scriptstyle *}$}}}

\newcommand{\Ran}{\operatorname{Ran}}
\DeclareMathOperator{\carr}{carr}
\DeclareMathOperator{\supp}{supp}
\DeclareMathOperator{\esupp}{ess-supp}

\DeclareMathOperator{\ad}{ad}
\DeclareMathOperator{\Inn}{Inn}
\DeclareMathOperator{\Tr}{Tr}
\DeclareMathOperator{\sgn}{sgn}
\DeclareMathOperator{\sd}{sd}
\DeclareMathOperator{\pr}{pr}

\DeclareMathOperator{\mindist}{min-dist}
\DeclareMathOperator{\maxdist}{max-dist}
\DeclareMathOperator{\minmaxdist}{\stack{\text{min}}{\text{max}}-dist}
\DeclareMathOperator{\maxmindist}{\stack{\text{max}}{\text{min}}-dist}

\renewcommand{\Re}{\operatorname{Re}}

\newcommand{\ip}[2]{{\big\langle #1 \, \big|\, #2\big\rangle}}
\newcommand{\ket}[1]{{\vert #1\rangle}}
\newcommand{\bra}[1]{{\langle #1 \vert}}
\newcommand{\stack}[2]{\substack{#1 \\ #2}}
\newcommand{\expt}[2]{{\langle #1 \rangle}_{#2}}

\newcommand{\ub}{\overline{u}}
\newcommand{\vb}{\overline{v}}
\newcommand{\wb}{\overline{w}}
%\newcommand{\Bb}{\overline{B}}
%\newcommand{\Pb}{\overline{P}}

%\newcommand{\Qb}{\overline{Q}}
\newcommand{\Xb}{\overline{X}}
\newcommand{\Xib}{\overline{\Xi}}
\newcommand{\Xxb}{\overline{\Xx}}
\newcommand{\Yb}{\overline{Y}}

\newcommand{\fhat}{\widehat{f}}
\newcommand{\hpsi}{\hat{\psi}}

\newcommand{\Ob}{{\boldsymbol{0}}}
\newcommand{\eb}{{\boldsymbol{e}}}
\newcommand{\Bb}{{\boldsymbol{B}}}
\newcommand{\Cb}{{\boldsymbol{C}}}
\newcommand{\Db}{{\boldsymbol{D}}}
\newcommand{\Fb}{{\boldsymbol{F}}}
\newcommand{\Ib}{{\boldsymbol{I}}}
\newcommand{\Lb}{{\boldsymbol{L}}}
\newcommand{\Mb}{{\boldsymbol{M}}}
\newcommand{\Nb}{{\boldsymbol{N}}}
\newcommand{\Pb}{{\boldsymbol{P}}}
\newcommand{\Qb}{{\boldsymbol{Q}}}
\newcommand{\Ub}{{\boldsymbol{U}}}


\newcommand{\Mbu}{{\boldsymbol{\utilde{M}}}}
\newcommand{\Nbu}{{\boldsymbol{\utilde{N}}}}

\newcommand{\Lc}{{\mathcal{L}}}
\newcommand{\Mc}{{\mathcal{M}}}
\newcommand{\Nc}{{\mathcal{N}}}
\newcommand{\Sc}{{\mathcal{S}}}
\newcommand{\Wc}{{\mathcal{W}}}

\newcommand{\Tu}{{\mathcal T}}

% Categories
\newcommand{\Bund}{{\sf Bund}}
\newcommand{\Ct}{{\sf C}}
\newcommand{\Jt}{{\sf J}}
\newcommand{\It}{{\sf I}}
\newcommand{\Ut}{{\sf U}}
\newcommand{\Cat}{{\sf Cat}}
\newcommand{\FSetInj}{{\sf FSet\hbox{-}Inj}}


\newcommand{\BkGrnd}{{\sf BkGrnd}}
\newcommand{\LCTo}{{\sf LCT}_0}
\newcommand{\LCT}{{\sf LCT}}
\newcommand{\Loco}{{\sf Loc}_0}
\newcommand{\Loc}{{\sf Loc}}
\newcommand{\CLoc}{{\sf CLoc}}
\newcommand{\FLoc}{{\sf FLoc}}
\newcommand{\SpinLoc}{{\sf SpinLoc}}
%\newcommand{\Loc}{\Loc}
%\newcommand{\Loco}{{\sf Man}}
%\newcommand{\Loc}{{\sf Man}^\sqcup}
\newcommand{\Obs}{{\sf Obs}}
\newcommand{\Val}{{\sf Val}}
\newcommand{\IsoP}{{\sf IsoPhys}}
%\newcommand{\Sympl}{{\sf CSympl}}
\newcommand{\preSympl}{{\sf preSympl}}

\newcommand{\Grp}{{\sf Grp}}
\newcommand{\Set}{{\sf Set}}
\newcommand{\iSet}{{\sf iSet}}
\newcommand{\States}{{\sf States}}
\newcommand{\Sympl}{{\sf Sympl}}
\newcommand{\Sys}{{\sf Sys}}
\newcommand{\SysLPE}{{\sf SysLPE}}
\newcommand{\Alg}{{\sf Alg}}
\newcommand{\CAlg}{{\sf C^*\hbox{-}Alg}}
\newcommand{\TAlg}{{\sf TAlg}}
\newcommand{\Test}{{\sf Test}}
\newcommand{\Top}{{\sf Top}}
\newcommand{\TVS}{{\sf TopVS}}
\newcommand{\GH}{{\sf GH}}
\newcommand{\Vect}{{\sf Vect}}
\newcommand{\iVect}{{\sf iVect}}
\newcommand{\Phys}{{\sf Phys}}
\newcommand{\twAlgSts}{{\sf gr\hbox{-}Alg\hbox{-}Sts}}
\newcommand{\xLCT}{{\sf xLCT}}
\newcommand{\xAlg}{{\sf xAlg}}
% Functors

\newcommand{\Af}{{\mathscr A}}
\newcommand{\Afd}{{\mathscr A}^\sqcup}
\newcommand{\Bf}{{\mathscr B}}
\newcommand{\Bfd}{{\mathscr B}^\sqcup}
\newcommand{\Cf}{{\mathscr C}}
\newcommand{\Df}{{\mathscr D}}
\newcommand{\Ef}{{\mathscr E}}
\newcommand{\Ff}{{\mathscr F}}
\newcommand{\Gf}{{\mathscr G}}
\newcommand{\If}{{\mathscr I}}
\newcommand{\Jf}{{\mathscr J}}
\newcommand{\Lf}{{\mathscr L}}
\newcommand{\Mf}{{\mathscr M}}
\newcommand{\Pf}{{\mathscr P}}
\newcommand{\Qf}{{\mathscr Q}}
\newcommand{\Rf}{{\mathscr R}}
\newcommand{\Sf}{{\mathscr S}}
\newcommand{\Tf}{{\mathscr T}}
\newcommand{\Uf}{{\mathscr U}}
\newcommand{\Vf}{{\mathscr V}}

\newcommand{\FfL}{\Ff_{\sf L}}
\newcommand{\FfS}{\Ff_{\sf S}}


\newcommand{\Rc}{{\sf R}}
\newcommand{\Tc}{{\mathcal T}}
\newcommand{\Vc}{{\mathcal V}}
\newcommand{\Wf}{{\mathscr W}}
\newcommand{\Zf}{{\mathscr Z}}

\newcommand{\zetad}{\zeta^\sqcup}
\newcommand{\etad}{\eta^\sqcup}

\newcommand{\CPT}{\mathscr{CPT}}
\newcommand{\CoPow}{\text{\bf CoPow}}

\newcommand{\PT}{\mathscr{PT}}
\DeclareMathOperator{\Sol}{Sol}
%\newcommand{\Sol}{{\mathscr L}}

\newcommand{\Forget}{{\mathscr V}}

\newcommand{\id}{{\rm id}}

\DeclareMathOperator{\Aut}{Aut}
\DeclareMathOperator{\Out}{Out}

\DeclareMathOperator{\Char}{Char}
\DeclareMathOperator{\Ext}{Ext}
\DeclareMathOperator{\Diff}{Diff}
\DeclareMathOperator{\Fol}{Fol}
\DeclareMathOperator{\Aff}{Aff}
\DeclareMathOperator{\cl}{cl}
\DeclareMathOperator{\opcl}{op-cl}
\DeclareMathOperator{\lincl}{lin-cl}
\DeclareMathOperator{\co}{co}
\DeclareMathOperator{\eq}{eq}
\DeclareMathOperator{\inte}{int}
\newcommand{\rce}{{\rm rce}}
\newcommand{\rced}{{\rm rce}^\sqcup}
\newcommand{\loc}{{\rm loc}}

\newcommand{\tc}{{\mathrm{tc}}}

\DeclareMathOperator{\dcl}{-cl}
\newcommand{\wscl}{{\rm w}^*\!\dcl}
%\newcommand{\wscl}{{\rm w}^*\textrm{-cl}}
\newcommand{\septens}{\stackrel{{\rm sep}}{\otimes}}

\newcommand{\cPhi}{\check{\Phi}}

\newcommand{\tchi}{\widetilde{\chi}}

\makeatletter
\newcommand{\raisemath}[1]{\mathpalette{\raisem@th{#1}}}
\newcommand{\raisem@th}[3]{\raisebox{#1}{$#2#3$}}
\makeatother
\newcommand{\act}[2]{\leftidx{^{\raisemath{-1.5pt}{#1}}}{#2}{}}
\newcommand{\adj}[1]{\leftidx{^{t}}{#1}{}}

\newcommand{\GL}{{\rm GL}}
\newcommand{\SL}{{\rm SL}}
\newcommand{\SU}{{\rm SU}}
\newcommand{\SO}{{\rm SO}}
\newcommand{\Spin}{{\rm Spin}}

\DeclareMathOperator{\lcm}{lcm}

\newcommand{\psum}{\sideset{}{'}\sum}
\DeclareMathOperator{\CCR}{CCR}
\newcommand{\WW}{\mathcal{W}}

\newcommand{\zett}{\widetilde{\zeta}}
\newcommand{\ett}{\widetilde{\eta}}
\newcommand{\xit}{\widetilde{\xi}}

\newcommand{\phit}{\widetilde{\phi}}


\newcommand{\cPsi}{\check{\Psi}}
\newcommand{\tS}{\widetilde{S}}
\newcommand{\elb}{{\boldsymbol{\ell}}}

\newcommand{\tpsi}{\widetilde{\psi}}
\newcommand{\tTf}{\widetilde{\Tf}}

\newcommand{\tJ}{\tilde{J}}
\renewcommand{\theenumi}{(\alph{enumi})}

\newcommand{\dd}{\mathrm{d}}
\DeclareMathOperator{\Ind}{Index}

\newcommand{\Ps}{\mathcal{P}}

\begin{document}
\begin{center}\Large Holomorphic compactly supported regularizing perturbations of Green-hyperbolic operators with discrete exceptional parameters\end{center}
\noindent\textbf{Stable ID:} 1148\quad\textbf{Author:} Christopher J. Fewster\par
\noindent\textbf{Work:} Modified Green-Hyperbolic Operators\par
\noindent\textbf{Source:} \url{https://sigma-journal.com/2023/057/}\par
\noindent\textbf{Version:} Official SIGMA 19 (2023) article 057, DOI 10.3842/SIGMA.2023.057; journal PDF and native TeX acquired 2026-09-30, exact bytes pinned by SHA256\par
\noindent\textbf{Rights:} CC BY 4.0. Authors retain copyright. \url{https://creativecommons.org/licenses/by/4.0/}. Reuse and typesetting adaptation; original language and mathematical content preserved.\par
\noindent\textbf{Difficulty (prior selection preserved):} {\CJKfont H3: The proof synchronizes analytic Fredholm inverses across every Sobolev order and upgrades them to the bounded-convergence topology of smooth compactly supported functions. It extends the compact-support inverse to the LF test-function space, constructs both Green operators and proves the precise causal support alternative. A separate holomorphic support argument recovers ordinary Green propagation for support-nonincreasing modifications, so this is more than applying the analytic Fredholm theorem at one Hilbert-space level.}\par
\medskip\noindent\textbf{Editorial boundary note (not author text):} Complete original principal Theorem3.1, all assumptions(a)–(d), exact delta=beta+gamma>0 and Sobolev thresholds, every conclusion(A)–(E), and the stronger support-nonincreasing(d-prime) alternative. Full globally hyperbolic/finitely connected spacetime, compact topologically regular K, causal/support/dual-density and Green-identity conventions, K-nonlocal definition, current Sobolev inversion Theorem5.1, smooth inversion Corollary5.2 and LF inverse-transfer Lemma5.3 with complete proofs are included. The entire six-step main proof preserves smooth-kernel finiteness, exceptional-parameter sets, smooth/LF holomorphic inversion, both Green identities, exact causal-support alternatives, continuous Sobolev extensions and ordinary-propagation argument. Complete source AppendixA definitions and every bounded-convergence/Sobolev-to-smooth/LF/Leibniz proof are retained. Green-hyperbolic extension and analytic Fredholm theory remain exact author-stated established prerequisites. One expressly named main perturbation outcome; later dual-kernel equality, integral-kernel examples and LU/application claims are unselected, with supporting lemmas receiving no extra counts. Original author support claims, spelling, inclusions and signs are preserved without mathematical repair.\par
\medskip\hrule\medskip\noindent\textbf{Author text begins below.}\par
\setcounter{section}{1}
% Original source lines 493--546
 \section{Green-hyperbolic operators}\label{sec:GHOs}

{\bf Preliminaries.}
 We begin by recalling the general setting of Green-hyperbolic operators~\cite{Baer:2015}. Let~$M$ be a smooth finite-dimensional manifold, allowing the possibility that $M$ has finitely many connected components with possibly different dimensions, and let $g$ be a smooth Lorentzian metric on $M$ of signature $+--\cdots$. With these structures, $M$ is automatically Hausdorff and paracompact~\cite{Geroch_spinorI:1968}.
 We assume that $(M,g)$ is time-orientable and that a time-orientation has been chosen. To minimise notation, we denote the Lorentzian spacetime formed by the manifold, metric and time orientation with the single symbol $M$. The volume measure induced by the metric will be denoted $\mu$. On other points of notation,
 the symbol $\subset$ will always allow for the possibility of equality, while $\NN_0$ and $\NN$ denote the natural numbers with or without zero, respectively.

 As usual, the causal future/past of a point $x\in M$ is denoted $J^\pm(x)$ and comprises all points (including $x$) that may be reached from $x$ along smooth future/past-directed curves. (Throughout this paper, we tacitly order alternatives labelled by $\pm$ or $\mp$ so that the alternative labelled by the upper symbol comes first.)
 If $S\subset M$ then one writes $J^\pm(S)=\cup_{x\in S} J^\pm(x)$, and $J(S)=J^+(S)\cup J^-(S)$. The spacetime is globally hyperbolic if it contains no causal curves and $J^+(K)\cap J^-(K)$ is compact for all compact sets $K\subset M$~\cite{Bernal:2006xf}. Globally hyperbolic spacetimes can be foliated into smooth spacelike Cauchy surfaces~\cite{Bernal:2004gm}; it is also the case that $J^\pm(K)$ are closed whenever $K$ is compact. We adopt the following terminology: $S\subset M$ is
 \emph{spacelike compact} if $S$ is closed and $S\subset J(K)$ for some compact $K$; $S$ is \emph{future/past-compact} if $S\cap J^{\pm}(x)$ is compact for all $x\in M$; $S$ is \emph{strictly future/past-compact} if
 $S\subset J^{\mp}(K)$ for some compact $K$.

 If $B$ is a vector bundle with finite-dimensional fibres, then $\Gamma^\infty(B)$ will denote the corresponding space of smooth sections and $\Gamma_{0/\pc/\fc/\spc/\sfc/\mathrm{sc}}^\infty(B)$ will be the spaces of smooth sections with com\-pact/past-compact/future-compact/strictly past-compact/strictly future-compact/space\-li\-ke-compact support. The space of smooth sections with support contained in some a closed subset $A\subset M$ is denoted $\Gamma^\infty_A(M)$. Further details on the topologies of these spaces are summarised in Appendix~\ref{appx:topspaces}. We will only consider bundles with finite-dimensional fibres that are (without any real loss) vector spaces over $\CC$. The bilinear (not sesquilinear!) pairing between sections of dual bundle $B^*$ and sections of $B$ is denoted with angle brackets $\langle \cdot,\cdot\rangle \colon \Gamma^\infty(B^*)\times \Gamma^\infty(B)\to C^\infty(M)$.

{\bf Green-hyperbolicity.} Suppose that $B_1$ and $B_2$ are bundles over $M$ and that $P\colon \Gamma^\infty(B_1)\to \Gamma^\infty(B_2)$ is a linear partial differential operator. Then there is a formal dual $\adj{P}\colon \Gamma^\infty(B_2^*)\to \Gamma^\infty(B_1^*)$ given by
 \begin{equation*}
 	\int_M \mu\, \big\langle \adj{P}f, \phi\big\rangle = \int_M\mu\, \langle f, P\phi\rangle
 \end{equation*}
 for all $\phi\in \Gamma^\infty(B_1)$ and $f\in \Gamma^\infty(B_2^*)$ whose supports intersect compactly; $\adj{P}$ is also a linear partial differential operator.
 If they exist, linear maps $E^\pm\colon \GoinX{B_2}\to \Gamma^\infty(B_1)$ obeying
 \begin{itemize}\itemsep=0pt
 	\item[(G1)] $E^\pm P f=f$ for all $f\in\GoinX{B_1}$,
 	\item[(G2)] $P E^\pm f = f$ for all $f\in\GoinX{B_2}$,
 	\item[(G3)] $\supp E^\pm f \subset J^\pm (\supp f)$ for all $f\in\GoinX{B_2}$
 \end{itemize}
 are called advanced ($-$) and retarded ($+$) Green operators for $P$.\footnote{B\"ar reverses the usage of `advanced' and `retarded'; we adopt the more standard convention.}
 \begin{Definition}
 	$P$ is said to be \emph{Green-hyperbolic} if both $P$ and $\adj{P}$ admit advanced and retarded Green operators.
 \end{Definition}

 It is a remarkable fact that this purely algebraic definition -- only requiring linearity of $E^\pm$ and with no presumption of uniqueness or continuity -- has the following consequence, which is a summary of results in~\cite[Sections 3 and 4]{Baer:2015} and deserves to be called the \emph{first main theorem of Green-hyperbolicity}.
 \begin{Theorem}\label{thm:GHbasics}
 	Let $P\colon\Gamma^\infty(B_1)\to \Gamma^\infty(B_2)$ be a Green-hyperbolic operator and consider:
 	\begin{enumerate}\renewcommand{\theenumi}{(\roman{enumi})}\itemsep=0pt
 		\item[$(i)$] the restriction $P\colon\Gamma^\infty_{\spc/\sfc}(B_1)\to\Gamma^\infty_{\spc/\sfc}(B_2)$,
 		\item[$(ii)$] the restriction $P\colon\Gamma^\infty_{\pc/\fc}(B_1)\to\Gamma^\infty_{\pc/\fc}(B_2)$, or
 		\item[$(iii)$] the extension $P\colon\DD'_{\pc/fc}(B_1)\to\DD'_{\pc/fc}(B_2)$ $($see below$)$.
 	\end{enumerate}
 	Then, in each case, $P$ has continuous inverses, denoted
 	$\tilde{E}^\pm$, $\overline{E}^\pm$ and $\widehat{E}^\pm$ in cases
 	$(i)$, $(ii)$, $(iii)$ respectively, and which are successive extensions of $E^\pm$. Each of these inverses has the support property $(G3)$, replacing
 	$\Gamma_0^\infty(B_2)$ by the appropriate domain. In particular,
 	$E^\pm$ are uniquely determined by $P$ and are continuous.
 \end{Theorem}	
 In $(iii)$ the space of distributional sections $\DD'(B)$ of a bundle $B$ over $M$ is defined
 as the topological dual of $\Gamma_0^\infty(B^*\otimes\Omega)$, where $\Omega$ is the bundle of weight-$1$ densities over $M$. As sections of $B$ and $B^*\otimes\Omega$ can be paired to give a density, there is an obvious embedding of $\Gamma^\infty(B)$ in $\DD'(B)$, with $\phi\in\Gamma^\infty(B)$ corresponding to the distribution
 $f\mapsto \int_M \langle \phi,f\rangle$ acting on $f\in \Gamma_0^\infty(B^*\otimes\Omega)$.
 The map $P$ in $(iii)$ is the restriction of the dual map
 $\big(\mu(\adj{P})\mu^{-1}\big)'$ to elements of $\DD'(B_1)$ with past- or future-compact support.
 (Recall that the formal dual is defined relative to the specific density~$\mu$.)
 In~\cite{Baer:2015}, B\"ar defines distributional sections of $B$ using
 the topological dual of $\Gamma_0^\infty(B^*)$, which would be the space of distribution densities in our terminology. The metric density
 provides an isomorphism between the spaces and operators that he considers and those that we do. As we have in mind potential applications where more than one metric might be in use, we have elected not to make a fixed identification between distributions and distribution densities.


\setcounter{Theorem}{3}
% Original source lines 578--640
 {\bf Modified Green-hyperbolic operators.}
 Turning to the subject of this paper, suppose $P\colon \Gamma^\infty(B_1)\to \Gamma^\infty(B_2)$ is Green-hyperbolic, with Green operators $E^\pm$. Let $A\colon\Gamma^\infty(B_1)\to \Gamma^\infty(B_2)$ be linear, with range contained in $\Gamma^\infty_K(B_2)$ for some compact, topologically regular $K\subset M$.

 Although $P+A$ is not necessarily a differential operator, we wish to state conditions analogous to (G1)--(G3) that can characterise suitable Green operators.
 To gain some insight, let us assume for the moment that
 for each $h\in \GoinX{B_2}$, the equation $(P+A)\phi=h$ has unique solutions with past/future-compact support given by $\phi=E_{P+A}^\pm h$ where $E_{P+A}^{\pm}\colon \GoinX{B_2}\to \Gamma^\infty(B_1)$ are linear maps with ranges
 necessarily contained in $\Gamma^\infty_{pc/fc}(B_1)$. Then we may write
 \begin{equation*}
 	PE_{P+A}^\pm f = f- AE_{P+A}^\pm f = PE_P^\pm\big(f- AE_{P+A}^\pm f\big),\qquad f\in \GoinX{B_2},
 \end{equation*}
 using our assumptions on $P$ and $A$. As $P$ is invertible on $\Gamma^\infty_{\pc/\fc}(B_1)$ by Theorem~\ref{thm:GHbasics}\,$(ii)$, we have
 \begin{equation*}
 	E_{P+A}^\pm f = E_P^\pm \big(f-AE_{P+A}^\pm f\big)
 \end{equation*}
 with support determined using condition (G3) for $P$,
 \begin{equation*}
 	\supp E_{P+A}^\pm f \subset J^\pm \big(\supp \big(f-AE_{P+A}^\pm f\big)\big) \subset J^\pm(\supp f\cup K).
 \end{equation*}
 By imposing further conditions on $A$ we may refine the information available. Specifically, suppose that $\phi|_K\equiv 0$ implies that $A\phi\equiv 0$.
 Observe that
 \begin{equation*}
 	(P+A)E_P^\pm f = f+ AE_P^\pm f \in \GoinX{B_1}
 \end{equation*}
 by (G2) for $E_P^\pm$ and the definition of $A$; by assumption on $E^\pm_{P+A}$ we now have
 \begin{equation*}
 	E_P^\pm f = E_{P+A}^\pm \big(f+ AE_P^\pm f\big)
 \end{equation*}
 and deduce that $E_{P+A}^\pm f=E_P^\pm f$ for all $f\in \GoinX{B_2}$ such that
 $J^\pm(\supp f)\cap K$ is empty \big(because~$AE_P^\pm f$ vanishes\big). Together with our earlier observation, $E^\pm_{P+A}$ satisfy the modified support property
 \begin{itemize}\setlength{\leftskip}{0.10cm}
\item[(G3${}'$)] for all $f\in \GoinX{B_2}$,
 	\begin{equation*}\displaystyle \supp E_{P+A}^{\pm} f \subset
 		\begin{cases} J^{\pm}(\supp f), & J^\pm(\supp f)\cap K=\varnothing,\\
 			J^{\pm}(\supp f\cup K), & \text{otherwise.}
 	\end{cases}\end{equation*}
 \end{itemize}
 With this intuition established, we can drop the assumption that $(P+A)\phi=h$ has unique solutions with past/future-compact support. The standing assumptions are now
 \begin{itemize}
 	\item[(A1)] $P\colon\Gamma^\infty(B_1)\to \Gamma^\infty(B_2)$ is Green-hyperbolic,
 	\item[(A2)] $A\colon\Gamma^\infty(B_1)\to \Gamma^\infty(B_2)$ is linear, with range contained in $\Gamma^\infty_K(B_2)$,
 	\item[(A3)] for $\phi\in \Gamma^\infty(B_1)$, $\phi|_K\equiv 0$ implies $A\phi\equiv 0$.
 \end{itemize}
 We now make the following definition.
 \begin{Definition}\label{def:KnonlocalGops}
 	Subject to assumptions (A1)--(A3), linear maps $E_{P+A}^{\pm}\colon\GoinX{B_2}\to \Gamma^\infty(B_1)$ are said to be \emph{retarded/advanced $K$-nonlocal Green operators} for $P+A$ if they satisfy~(G1) and~(G2) \big(with $P+A$ replacing $P$ and $E_{P+A}^\pm$ replacing $E^\pm$\big) and (G3${}')$. If both $P+A$ and $\adj{P}+\adj{A}$ admit retarded and advanced $K$-nonlocal Green operators then $P+A$ is called \emph{$K$-nonlocally Green-hyperbolic}.
 \end{Definition}



 As with Green-hyperbolic operators, the above definition implies considerably more and indeed the analogue of Theorem~\ref{thm:GHbasics} holds with (G3) replaced by (G3${}')$. These results will be proved in~\cite{FewVer:20xx}. Our main goal here will be to give sufficient
 conditions for the existence of $K$-nonlocal Green operators. For simplicity of presentation, we restrict to operators acting on spaces of smooth functions, rather than bundle sections. We will also establish some continuity results for the $K$-nonlocal Green operators, both from $C_0^\infty(M)$ to $C^\infty(M)$ and between various Sobolev spaces, and address the
 holomorphicity of the $K$-nonlocal Green operators with respect to a~parameter.

 At a formal level it is straightforward to see what the Green operators of $P+A$ should be, if they exist: if
 $\phi\in C^\infty_{\pc/\fc}(M)$ solves $(P+A)\phi=f\in \CoinX{M}$, then
 \begin{equation*}
 	\phi = E_P^\pm (f- A\phi) = E^\pm_{P} g,
 \end{equation*}
 where $g=f - A \phi= f- AE_P^\pm g$, so \smash{$g=\big(I+AE_P^\pm\big)^{-1}f$}. Formally, therefore, the Green operators for $P+A$ are
 \begin{equation}\label{eq:formalGF}
 	E_{P+A}^\pm = E_P^\pm \big(I+AE^\pm_{P}\big)^{-1},
 \end{equation}
 and the technical task is to make this formula rigorous, where possible, and to establish that the resulting operators $E_{P+A}^\pm$ are indeed $K$-nonlocal Green operators for $P+A$. Note that the inversion of $I+AE^\pm$ must be performed in $\CoinX{M}$. We will accomplish this by first inverting in Sobolev spaces, where Hilbert space techniques can be used, and then boosting the result up to $\CoinX{M}$.


% Original source lines 644--701
 \section{Main result and remarks}\label{sec:main}

 Let $M$ be a globally hyperbolic spacetime with at most finitely many connected components,
 and let $K$ be a fixed, topologically regular, compact subset of $M$. For brevity, we write $C^\infty$ for~$C^\infty(M)$ and so on; for a closed set $A$, $C^\infty_A$ denotes the space of smooth functions with support in $A$, and the same convention is used for other spaces of distributions and Sobolev spaces. If $X$ and $Y$ are locally convex topological spaces, $\LL_b(X,Y)$ denotes the space of linear maps between them, equipped with the topology of bounded convergence (see Appendix~\ref{appx:topspaces}); for linear maps between normed spaces, this coincides with the usual operator norm topology. As usual, we write $\LL_b(X)$ for $\LL_b(X,X)$.

 The main result of this paper is the following.
 \begin{Theorem} \label{thm:GHpert}
 	Let $P$ be a Green-hyperbolic operator and
 	suppose $A(\lambda)$, $\lambda\in\CC$, is a family of linear maps $A(\lambda)\colon C^\infty\to C^\infty$ whose ranges are contained in $C^\infty_K$ and with $A(0)=0$. Suppose that,
 	for some $\beta,\gamma\in \RR$ with $\delta=\beta+\gamma>0$, and some $s_*\in \RR$,
 	\begin{enumerate}\itemsep=0pt
 		\item[$(a)$] the Green operators $E^\pm$ of $P$ extend to linear maps
 		$\EE'\to \DD'$ with continuous restrictions mapping $H^s_0\to H^{s+\beta}_\loc$ for all $s\ge s_*$;
 		\item[$(b)$] each map $A(\lambda)$ extends to a linear map $\DD'\to\DD'_K$ with continuous restrictions mapping $H^s_\loc\to H^{s+\gamma}_K$ for all $s\ge s_*+\beta$
 		$($consequently, the continuous maps $A(\lambda)E^\pm\colon H_0^s\to H_K^{s+\delta}$ induce
 		compact maps $A(\lambda)E^\pm\colon H_0^s\to H_K^{s}$ due to the Sobolev embedding theorems and $\delta>0)$;
 		\item[$(c)$] for all $s\ge s_*$, the compact map $H^s_0\to H^s_K$ induced by the compositions $A(\lambda)E^\pm$ is holomorphic in $\lambda\in\CC$ with respect to the topology of $\LL_b(H^s_0,H^s_K)$;
 		\item[$(d)$] if $f\in C^\infty$ vanishes identically on $K$ then $Af=0$.
 	\end{enumerate}
 	Then
 	\begin{enumerate}\renewcommand{\labelenumi}{(\Alph{enumi}).}
 		\item[$(A)$] $\ker (P+ A(\lambda))|_{C^\infty_{\pc/\fc}}$ has finite dimension
 		for each $\lambda\in\CC$;
 		\item[$(B)$]%\label{it:discset}
 the sets
 		\begin{equation*}
 			S^\pm = \big\{\lambda\in \CC\colon \ker (P+ A(\lambda))|_{C^\infty_{\pc/\fc}}\neq 0\big\} \qquad \text{and}\qquad
 			S = S^+\cup S^-
 		\end{equation*}
 		are discrete subsets of $\CC$, whose complements are open neighbourhoods of zero;
 		\item[$(C)$]%\label{it:Gops}
 $P+A(\lambda)$ has advanced and retarded $K$-nonlocal Green operators for all $\lambda\in\CC\setminus S$, that are holomorphic in $\lambda$ on this domain with respect to the topology of $\LL_b(C_0^\infty, C^\infty)$; these Green operators are given explicitly as
 		\begin{equation*}
 			\tilde{E}_{\lambda}^\pm = E^\pm \big(I+A(\lambda)E^\pm\big)^{-1},
 		\end{equation*}
 		where the inverse on the right-hand side exists and is holomorphic in $\lambda\in \CC\setminus S$
 		with respect to the topology of $\LL_b(C_0^\infty)$;
 		\item[$(D)$]%\label{it:Gopsext}
 the Green operators
 		possess continuous extensions mapping $H^s_0\to H^{s+\beta}_\loc$ for all $s\ge s_*$ and which are holomorphic in $\lambda\in\CC\setminus S$ with respect to the topology of $\LL_b\big(H^s_0,H^{s+\beta}_\loc\big)$;
 		
 		\item[$(E)$]%\label{it:locGops}
 if $(d)$ is replaced by
 		\begin{enumerate}
 			\item[$(d')$] $\supp A(\lambda)f\subset \supp f$ for all $f\in C^\infty$, and $\lambda\in\CC$
 		\end{enumerate}
 		then $P+ A(\lambda)$ has advanced and retarded Green operators in the usual sense
 		for $\lambda\in\CC\setminus S$.
 	\end{enumerate}
 \end{Theorem}
 Of course, in the situation where the hypotheses also apply to the formal duals of $P$ and $A(\lambda)$ then the upshot is that $P+ A(\lambda)$ and $\adj{P}+\adj{A}(\lambda)$
 both admit advanced and retarded $K$-nonlocal Green operators for all
 $\lambda\in\CC$ in an open $0$-neighbourhood with discrete complement given by the union of $S$ with the corresponding set for $\adj{P}$. For the non-exceptional values, $P+A(\lambda)$ is $K$-nonlocally Green-hyperbolic. The general theory of such operators and their properties is developed in detail in~\cite{FewVer:20xx}. We note in particular that when
 the ranges of both $A(\lambda)$ and $\adj{A}(\lambda)$ are contained in $C^\infty_K$ then assumption $(d)$ holds automatically for both operators.

 In many situations the holomorphicity assumption $(c)$ is easily verified, e.g., where $A(\lambda)$ is a~polynomial in $\lambda$. Precomposing with the continuous embedding $H^s_K\to H^s_0$, it follows that the restrictions of $A(\lambda)E^\pm$ to $H^s_K$, for $s\ge s_*$, are holomorphic with respect to the norm topology
 of $\LL(H^s_K)$. On the other hand, postcomposing with the embedding $H^s_K\to H^s_0$, we also have holomorphicity of $A(\lambda)E^\pm$ with respect to
 $\LL_b(H^s_0)$ (e.g., see Lemma~\ref{lem:Lb_basic}).

\setcounter{section}{4}
% Original source lines 922--1193
 \section{Proof of Theorem~\ref{thm:GHpert}}\label{sec:proof}

 The proof of Theorem~\ref{thm:GHpert} starts by establishing an inversion result that will later be applied to $I+A(\lambda)E^\pm$, where $E^\pm$ are the Green operators of $P$. This is proved first in Sobolev spaces $H^s_K$, using the analytic Fredholm theorem, and then extended to $C^\infty_K$. The main part of the proof then uses this information to define Green operators for $P+ A(\lambda)$.


 \begin{Theorem}\label{thm:Sobolevinversion}
 	Let $K\subset M$ be a fixed topologically regular compact set. Suppose
 	$Y(\lambda)$, $\lambda\in\CC$, is a~family of linear maps
 	$Y(\lambda)\colon\DD_K'\to\DD'_K$, each restricting to continuous maps $Y_s(\lambda)\colon H^s_K\to H^{s+\delta}_K$ for all $s\ge s_*\in\RR$ and some fixed $\delta>0$. Assume also that $Y(0)=0$.
 	Let $\hat{Y}_s(\lambda)\in\LL(H^s_K)$ denote the compact maps obtained by composing $Y_s(\lambda)$ with the embedding $H^{s+\delta}_K\hookrightarrow H^s_K$. Suppose that $\lambda \to \hat{Y}_s(\lambda)$ is holomorphic on $\CC$ in $\LL_b(H^s_K)$ for all $s\ge s_*$. Then
 	\begin{enumerate}\itemsep=0pt
 		\item[$(a)$] The set
 		\begin{equation}\label{eq:Sdef}
 			S=\big\{\lambda\in\CC\colon \ker (I+ Y(\lambda))|_{C^\infty_K}\neq 0\big\}
 		\end{equation}
 		is a discrete subset of $\CC$, and $\CC\setminus S$ is an open $0$-neighbourhood. Furthermore, $\ker (I+ Y(\lambda))|_{C^\infty_K}$ is finite-dimensional, and equal to
 		$\ker (I+ \hat{Y}_s(\lambda))$ for each $s\ge s_*$.
 		
 		\item[$(b)$] For all $\lambda\in \CC\setminus S$ and all $s\ge s_*$, the map $I+ \hat{Y}_{s}(\lambda)$ is continuously invertible for all $s\ge s_*$, and $\lambda\mapsto \big(I+ \hat{Y}_{s}(\lambda)\big)^{-1}$ is holomorphic in $\CC\setminus S$. Moreover,
 		$\big(I+ \hat{Y}_{s}(\lambda)\big)^{-1}$ is the restriction of $\big(I+ \hat{Y}_{s_*}(\lambda)\big)^{-1}$ to $H^s_K$ for each $s\ge s_*$, $\lambda\in\CC\setminus S$.
 		
 		\item[$(c)$] $(i)$ Suppose that $f\colon \CC\setminus S\to H^{s_*}_K$ is holomorphic and
 		there is a compact subset ${K_f\subset K}$ so that $\supp \big(\hat{Y}_{s_*}(\lambda)\big)^r f(\lambda)\subset K_f$ for all $r\in\NN_0$ and $\lambda\in \CC\setminus S$. Then $\supp \big(I+\hat{Y}_{s_*}(\lambda)\big)^{-1}f(\lambda)\subset K_f$ for all $\lambda\in \CC\setminus S$.
 		
 		$(ii)$ If $Y(\lambda)C^\infty_K\subset C^\infty_{K'}$ for some compact $K'\subset K$, then
 		$\supp \big(I+ \hat{Y}_{s_*}(\lambda)\big)^{-1}f\subset K'$ for all $\lambda\in \CC\setminus S$ and all $f\in C^\infty_K$, and
 		$\ker \big(I+ \hat{Y}_{s_*}(\lambda)\big)$ is a finite-dimensional subspace of
 		$C^\infty_{K'}$ for all $\lambda\in\CC$.
 		
 		NB. In $(c)$, `support' is to be understood as distributional support.
 	\end{enumerate}
 \end{Theorem}
 \begin{proof}
 	$(a)$ Compactness of $\hat{Y}_s(\lambda)$ follows from the Sobolev embedding theorems.
 	The function $\lambda\mapsto I+ \hat{Y}_{s_*}(\lambda)$ is an analytic function on $\CC$ with values in the Fredholm operators on $H^{s_*}_K$, and which is invertible for $\lambda=0$ because $\hat{Y}_{s_*}(0)=0$. By the analytic Fredholm theorem~\cite[Theorem VI.14]{ReedSimon:vol1}, $I+ \hat{Y}_{s_*}(\lambda)$ is invertible for all $\lambda\in\CC$ with the exception of the (possibly empty) set $S_{s_*}$ of $\lambda\in \CC$ for which $\ker \big(I+ \hat{Y}_{s_*}(\lambda)\big)$ is nontrivial (and necessarily finite-dimensional by~\cite[Theorem VI.15]{ReedSimon:vol1}). Furthermore, $S_{s_*}$ is a
 	discrete subset of $\CC$, whose complement is an open $0$-neighbourhood, and the inverse $\big(I+ \hat{Y}_{s_*}(\lambda)\big)^{-1}$ is meromorphic on $\CC$ and holomorphic on $\CC\setminus S_{s_*}$.
 	
 	If $f\in\ker \big(I+\hat{Y}_{s_*}(\lambda)\big)$ then $f = - \hat{Y}_{s_*}(\lambda)f\in H^{s_*+\delta}_K$ obeys $f\in\ker \big(I+ \hat{Y}_{s_*+\delta}(\lambda)\big)$; iterating, $f\in \bigcap_{s\ge s_*} H^s_K=C^\infty_K$ and $f\in \ker (I+ Y(\lambda))|_{C^\infty_K}$.
 	As the converse inclusion is trivial, we deduce that
 	$\ker \big(I+\hat{Y}_{s_*}(\lambda)\big)= \ker (I+ Y(\lambda))|_{C^\infty_K}$
 	for all $\lambda\in \CC$ and hence that $S_{s_*}=S$ defined in~\eqref{eq:Sdef}. In particular, for each fixed $\lambda\in\CC$, all the kernels $\ker \big(I+\hat{Y}_{s}(\lambda)\big)$ are equal for $s\ge s_*$.
 	
 	$(b)$ The first statement is immediate by the Fredholm theorem and part~(a); the second follows because $I+\hat{Y}_{s}(\lambda)$ coincides with $I+ \hat{Y}_{s_*}(\lambda)$ on $H^s_K$ and both operators are invertible for $\lambda\in\CC\setminus S$.
 	
 	$(c)$\,$(i)$ Let $\chi\in C_0^\infty$ vanish on $K_f$. Regarding each $\big(\hat{Y}_{s_*}(\lambda)\big)^{n}f(\lambda)$ as a distribution with support in $K_f$, we see (e.g., by \cite[Theorem 2.3.3]{Hormander1}) that
 	$\big(\hat{Y}_{s_*}(\lambda)^{n}f(\lambda)\big)(\chi)=0$ for all $n\in \NN_0$ and so the
 	holomorphic function on $\CC\setminus S$ defined by $\lambda\mapsto \big(\big(I+\hat{Y}_{s_*}(\lambda)\big)^{-1}f(\lambda)\big)(\chi)$ vanishes in a neighbourhood of the origin on which the resolvent may be expanded as a convergent geometric series in powers of $\hat{Y}_{s_*}(\lambda)$, recalling that $\hat{Y}_{s_*}(\lambda)\to 0$ as $\lambda\to 0$.
 	By holomorphicity, it follows that
 	$\big(\big(I+\hat{Y}_{s_*}(\lambda)\big)^{-1}f(\lambda)\big)(\chi)=0$ for all $\lambda\in\CC\setminus S$.
 	Allowing $\chi$ to vary, we conclude that $\supp \big(I+\hat{Y}_{s_*}(\lambda)\big)^{-1}f(\lambda)\subset K_f$.
 	
 	$(c)$\,$(ii)$ The assumption implies that of $(c)$\,$(i)$ for all $f\in C^\infty_K$ \big(regarded as constant functions from $\CC\setminus S$ to $H^{s_*}_K$\big) with $K_f=K'$, so
 	$\supp \big(I+\hat{Y}_{s_*}(\lambda)\big)^{-1}f\subset K'$ for all $\lambda\in \CC\setminus S$ and all $f\in C^\infty_K$. Finally,
 	if $g\in \ker \big(I+\hat{Y}_{s_*}(\lambda)\big)=\ker (I+ Y(\lambda))|_{C^\infty_K}$ (already known to be finite-dimensional) then $g\in C^\infty_K$ and $g=- Y(\lambda) g\in C^\infty_{K'}$.
 \end{proof}

 The hypotheses of Theorem~\ref{thm:Sobolevinversion} entail that each $Y(\lambda)$ restricts to a continuous endomorphism of $C^\infty_K$. To see this, note that for any $k\in\NN_0$ and any integer $s>k+n/2$, where $n$ is the maximum dimension of any component of $M$, there are continuous maps
 \begin{equation*}
 	C^s_K \longrightarrow H^s_K \stackrel{Y(\lambda)}{\longrightarrow} H^s_K \longrightarrow C^k_K,
 \end{equation*}
 where the unlabelled arrows are Sobolev embeddings. Thus for all $k\in \NN_0$ there is $s\in \NN_0$ and a constant $C$ such that
 $\|Y(\lambda)f\|_{K,k}\le C \|f\|_{K,s}$ for all $f\in C^\infty_K$. This observation allows the following conclusion to be drawn.
 \begin{Corollary}\label{cor:CKinfty_inversion}
 	In the notation of Theorem~$\ref{thm:Sobolevinversion}$, and for $\lambda\in\CC\setminus S$, $I+ Y(\lambda)$ restricts to a~homeomorphism of $C^\infty_K$, with inverse to be denoted $R(\lambda)$. The map $\lambda\mapsto R(\lambda)$ is
 	holomorphic on $\CC\setminus S$ in the topology of $\LL_b\big(C^\infty_K\big)$.
 	If, additionally, for some holomorphic function $f\colon \CC\setminus S\to C^\infty_K$
 	there is a compact subset $K_f\subset K$ so that $\supp Y(\lambda)^r f(\lambda)\subset K_f$ for all $r\in\NN_0$ and $\lambda\in \CC\setminus S$, then $\supp R(\lambda) f(\lambda)\subset K_f$ for all $\lambda\in \CC\setminus S$. If $Y(\lambda)C^\infty_K\subset C^\infty_{K'}$ for some compact $K'\subset K$, then
 	$\ker(I+ Y(\lambda))|_{C^\infty_K}$ is a finite-dimensional subspace of $C^\infty_{K'}$
 	for all $\lambda\in\CC$.
 \end{Corollary}
 \begin{proof}
 	Suppose $f\in C^\infty_K$ and $\lambda\notin S$. Then, for all $s\ge s_*$,
 	we have	$f\in H^s_K$ and $\big(I+ \hat{Y}_{s}(\lambda)\big)^{-1}f=\big(I+\hat{Y}_{s_*}(\lambda)\big)^{-1}f$, and consequently
 	\begin{equation*}
 		\big(I+\hat{Y}_{s_*}(\lambda)\big)^{-1}f\in \bigcap_{s\ge s_*} H^s_K = C^\infty_K.
 	\end{equation*}
 	Thus $\big(I+\hat{Y}_{s_*}(\lambda)\big)^{-1}$ restricts to linear self-map of $C^\infty_K$ which
 	\big(because $Y(\lambda)$ and $\hat{Y}_{s_*}(\lambda)$ agree on $C^\infty_K$\big) is a linear inverse to the restriction of $I+ Y(\lambda)$. Accordingly, $I+ Y(\lambda)$ restricts to a~continuous bijection of $C^\infty_K$, and since the latter is a Fr\'echet space the inverse mapping theorem implies that the inverse $R(\lambda)$ is continuous.
 	As every $C^\infty_K$ semi-norm is dominated by a~Sobolev norm and
 	vice versa, convergence in the operator norm of every $H^s_K$, for $s\ge s_*$,
 	implies convergence in $\LL_b\big(C^\infty_K\big)$, by Corollary~\ref{cor:HstoCinfty}. It follows that $\lambda\mapsto R(\lambda)$ is holomorphic on~$\CC\setminus S$ in the topology of $\LL_b\big(C^\infty_K\big)$.
 	The remaining statements are immediate from parts~$(a)$ and $(c)$ of Theorem~\ref{thm:Sobolevinversion} and the fact that distributional support of a smooth function is exactly its usual support.
 \end{proof}


 We need the following elementary observation, which will be used for $F=C_K^\infty$, $G=C_0^\infty$.
 \begin{Lemma}\label{lem:inversetransfer}
 	Let $F$ and $G$ be topological vector spaces and suppose that the diagram
 	\begin{equation*} %\label{eq:lemmacd}
 		\begin{tikzcd}
 			F \arrow[r, "S"]\arrow[d,"\iota"] & F \arrow[d,"\iota"]\\
 			G \arrow[r, "T"]\arrow[ur,"\hat{T}"] & G
 		\end{tikzcd}
 	\end{equation*}
 	of continuous linear maps commutes. If $\id_F+S$ is continuously invertible, then
 	$\id_G+T$ is continuously invertible with inverse
 	\begin{equation*}%\label{eq:invlemform}
 		(\id_G+ T)^{-1} = \id_G-T+\iota S(\id_F+S)^{-1}\hat{T}.
 	\end{equation*}
 \end{Lemma}
 \begin{proof}
 	Assuming that $\id_F+S$ is continuously invertible, we compute
 	\begin{align*}
 		(\id_G+T)\big(\id_G-T+\iota S(\id_F+S)^{-1}\hat{T}\big) = \id_G-T^2 + \iota (\id_F+S)S(\id_F+S)^{-1}\hat{T}=\id_G
 	\end{align*}
 	and
 	\begin{align*}
 		\big(\id_G-T+\iota S(\id_F+S)^{-1}\hat{T}\big)(\id_G+T) &= \id_G-T^2 + \iota S(\id_F+S)^{-1}(\id_F+S)\hat{T}= \id_G.
 	\end{align*}
 	using in both cases the identities $\iota S=T\iota$ and hence $\iota S \hat{T} = T\iota\hat{T}=T^2$, and also $\hat{T}T = \hat{T}\iota\hat{T} = S\hat{T}$. Therefore $\id_G+T$ has a linear inverse, given by a manifestly continuous expression.
 \end{proof}

 We come to the proof of the main result.
 \begin{proof}[Proof of Theorem~\ref{thm:GHpert}] The proof involves several steps, and uses some technical lemmas from the Appendix.
 	
 	\emph{$1.$ Preliminary observations.} Define $Y^\pm(\lambda)\in\LL(C_0^\infty,C^\infty)$ by
 	\begin{equation*}
 		Y^\pm(\lambda)=A(\lambda)E^\pm.
 	\end{equation*}
 	Owing to assumptions $(a)$ and $(b)$, we may extend $Y^\pm(\lambda)$ to linear maps $\EE'\to\DD'_K$ with continuous restrictions mapping $H^s_0\to H^{s+\delta}_K$ for all $s\ge s_*$; hence there are also continuous restrictions $Y_s^\pm(\lambda)\colon H^s_K\to H^{s+\delta}_K$ for $s\ge s_*$,
 	and by the Sobolev embedding theorems, compact maps $\hat{Y}^\pm_s(\lambda)\colon H^s_K\to H^{s}_K$ for $s\ge s_*$ which are holomorphic with respect to the operator norm topology, as noted after the statement of Theorem~\ref{thm:GHpert}.
 	
 	In fact, if $K'$ is any compact topologically regular set, the same argument shows that~$Y^\pm(\lambda)$
 define compact maps depending holomorphically on $\lambda$ in the topology of
 	$\LL_b\big(H^s_{K\cup K'}\big)$ for all~${s\!\ge\!s_*}$. By Corollary~\ref{cor:HstoCinfty}, it follows that
 	$Y^\pm(\lambda)$ are holomorphic in $\lambda$ with respect to the $\LL_b\big(C^\infty_{K\cup K'}\big)$ topology.
 	As $C^\infty_K$ and $C^\infty_{K'}$ are continuously embedded topological subspaces of $C^\infty_{K\cup K'}$, and $\Ran Y^\pm(\lambda)\subset C^\infty_K$, we may use Lemma~\ref{lem:Lb_basic}\,$(a)$,\,$(c)$ to deduce that
 	the $Y^\pm(\lambda)$ are holomorphic in $\lambda$ with respect to the
 	topology of $\LL_b\big(C^\infty_{K'},C^\infty_{K}\big)$.
 	By Lemma~\ref{lem:LbLF}, it follows that
 	$A(\lambda)E^\pm$ are also holomorphic with respect to $\LL_b\big(C_0^\infty,C^\infty_K\big)$ and $\LL_b(C_0^\infty)$; using Lemma~\ref{lem:Lb_basic}\,$(a)$, we also have holomorphicity with respect to
 	$\LL_b\big(C^\infty_K,C_0^\infty\big)$.
 	
 	\emph{$2.$ Finite-dimensionality of $\ker (P+ A(\lambda))|_{C^\infty_{\pc/\fc}}$.} Next, observe that $P$ induces bijections between $\ker (P+ A(\lambda))|_{C^\infty_{\pc/\fc}}$ and $\ker \big(I+ Y^{\pm}(\lambda)\big)\big|_{C^\infty_K}$,
 	with inverses given by the restrictions of~$E^{\pm}$. For
 	if $(P+ A(\lambda))\phi=0$ with $\phi\in C^\infty_{\pc/\fc}$ then
 	$P\phi=- A(\lambda)\phi\in C^\infty_K$ and $\phi=E^\pm P\phi$, so $P\phi\in\ker\big(I+ A(\lambda)E^\pm\big)\big|_{C^\infty_K}$; conversely, if $\big(I+ A(\lambda)E^\pm\big)h=0$ with
 	$h\in C^\infty_K$, then $PE^\pm h = h=- A(\lambda)E^\pm h$, so
 	$E^\pm h\in \ker(P+ A(\lambda))|_{C^\infty_{\pc/\fc}}$.
 	Thus
 	\begin{equation*}
 		S^\pm:= \big\{\lambda\in \CC\colon \ker (P+ A(\lambda))|_{C^\infty_{\pc/\fc}}\neq 0\big\}= \big\{\lambda\in \CC\colon \ker \big(I+ Y^{\pm}(\lambda)\big)\big|_{C^\infty_K}\neq 0\big\}.
 	\end{equation*}
 	In combination with Theorem~\ref{thm:Sobolevinversion}\,$(a)$, we also have
 	\begin{equation}\label{eq:GHpert_remark}
 		\dim \ker (P+A(\lambda))|_{C^\infty_{\pc/\fc}} = \dim \ker \big(I+Y_s^{\pm}(\lambda)\big)\big|_{C^\infty_{K}} = \dim \ker \big(I+\hat{Y}_s^{\pm}(\lambda)\big)<\infty
 	\end{equation}
 	\big(the latter kernel taken in $H^s_K$\big) for all $s\ge s_*$. Part~$(A)$ of the Theorem is thus proved.
 	
 	\emph{$3.$ Construction of holomorphic candidate $K$-nonlocal Green operators.} Applying Theorem~\ref{thm:Sobolevinversion} and Corollary~\ref{cor:CKinfty_inversion} to $Y^\pm(\lambda)$, one finds that $S^\pm$ are discrete subsets of $\CC$, whose complements in~$\CC$ are open $0$-neighbourhoods, and the operators $I+ Y^\pm(\lambda)$ are continuously invertible on~$C^\infty_K$ for $\lambda\in\CC\setminus S^\pm$, with inverses
 	that are holomorphic on~$\CC\setminus S^\pm$ in the topology
 	of $\LL_b\big(C^\infty_K\big)$. All these properties hold also for $S=S^+\cup S^-$,
 	so part~(B) is proved.
 	
 	Fixing any $\lambda\in \CC\setminus S$, the diagram of continuous maps
 	\begin{equation*}
 		\begin{tikzcd}
 			C^\infty_K \arrow[r, "T_K"]\arrow[d] & C^\infty_K \arrow[d]\\
 			C_0^\infty \arrow[r, "T_0"]\arrow[ur,"\hat{T}_0"] & C_0^\infty
 		\end{tikzcd}
 	\end{equation*}
 	commutes, where the unlabelled arrows are the canonical inclusions,
 	$T_0$ and $T_K$ are the respective restrictions of $Y^\pm(\lambda)$
 	to $C_0^\infty$, and $C_K^\infty$ and $\hat{T}_0$ exists because
 	$A(\lambda)C^\infty\subset C^\infty_K$.
 	It follows from Lemma~\ref{lem:inversetransfer} that $I+ Y^\pm(\lambda)$ are continuously invertible on $C_0^\infty$. By abuse of notation we write the inverses as $\big(I+ A(\lambda)E^\pm\big)^{-1}$; Lemma~\ref{lem:inversetransfer} now gives the identity
 	\begin{equation}\label{eq:inverseidentity}
 		\big(I+ A(\lambda)E^\pm\big)^{-1} = I - A(\lambda)E^\pm + A(\lambda)E^\pm \big(I+ A(\lambda)E^\pm\big)^{-1} A(\lambda)E^\pm,
 	\end{equation}
 	on $C_0^\infty(K)$, where we suppress notation for inclusions and restrictions. Note that the inverse on the right-hand side is taken in $\LL\big(C^\infty_K\big)$, while the left-hand side is an inverse in $\LL(C^\infty_0)$. Because the former inverse is holomorphic in
 	$\LL_b\big(C^\infty_K\big)$, the Leibniz rule (see Corollary~\ref{cor:Leibniz}) implies that the left-hand side is
 	holomorphic in $\LL_b(C_0^\infty)$; here, we have also used the holomorphicity
 	of~$A(\lambda)E^\pm$ in $\LL_b\big(C_0^\infty,C^\infty_K\big)$, $\LL_b(C_0^\infty)$ and $\LL_b\big(C^\infty_K,C_0^\infty\big)$ established in step 1 of the proof. It follows that the operators
 	\begin{equation*}
 		\tilde{E}^\pm_\lambda= E^\pm \big(I+ A(\lambda)E^\pm\big)^{-1}\in\LL(C_0^\infty, C^\infty),
 	\end{equation*}
 	which are the candidate Green operators for $P+A(\lambda)$,	are holomorphic in $\lambda$ on $\CC\setminus S$ with respect to the topology of $\LL_b(C_0^\infty,C^\infty)$.
 	To prove part (C) it is now enough to check that $\tilde{E}^\pm_\lambda$ are indeed $K$-nonlocal Green operators.
 	
 	\emph{$4.$ Verification that $\tilde{E}^\pm_\lambda$ are $K$-nonlocal Green operators.}
 	Given any $f\in C_0^\infty$,
 	\begin{equation*}
 		g=\big(I+ A(\lambda)E^\pm\big)^{-1} f
 	\end{equation*}
 	is the unique element of $C_0^\infty$ obeying
 	\begin{equation}\label{eq:geqn}
 		g + A(\lambda) E^\pm g= f,
 	\end{equation}
 	whereupon we deduce that $\supp g\subset \supp f \cup K$ and that
 	\begin{equation*}
 		\varphi = E^\pm g = E^\pm \big(I+ A(\lambda)E^\pm\big)^{-1} f = \tilde{E}^\pm_\lambda f
 	\end{equation*}
 	satisfies
 	\begin{equation}\label{eq:Pphi}
 		P\varphi + A(\lambda)\varphi = g+A(\lambda)E^\pm g=f, \qquad \supp \varphi\subset J^\pm(\supp g) \subset J^\pm(K\cup \supp f).
 	\end{equation}
 	In the special case $f=(P+A(\lambda))h$, for $h\in C_0^\infty$,
 	we have
 	\begin{equation*}
 		f = (P+ A(\lambda))E^\pm P h=Ph + A(\lambda) E^\pm Ph
 	\end{equation*}
 	and the unique solution to \eqref{eq:geqn} is clearly $g=Ph$. Thus
 	\begin{equation*}
 		Ph=g=\big(I+ A(\lambda)E^\pm\big)^{-1}f =\big(I+ A(\lambda)E^\pm\big)^{-1}(P+ A(\lambda))h;
 	\end{equation*}
 	consequently $E^\pm \big(I+ A(\lambda)E^\pm\big)^{-1} (P+ A(\lambda))h = h$.
 	In combination with~\eqref{eq:Pphi}, we have shown
 	\begin{equation*}
 		(P+ A(\lambda))\tilde{E}^\pm_\lambda f=f = \tilde{E}^\pm_\lambda (P+ A(\lambda)) f,\qquad\text{and}\qquad
 		\supp\tilde{E}^\pm_\lambda f\subset J^\pm(K\cup \supp f)
 	\end{equation*}
 	for all $f\in C_0^\infty$.
 	Now suppose more specifically that $J^\pm(\supp f)\cap K=\varnothing$ for $f\in C^\infty_0$. Then $A(\lambda)E^\pm f=0$, due to assumption $(d)$, and~\eqref{eq:geqn} is solved by $g=f$, so
 	$\tilde{E}^\pm_\lambda f=E^\pm f$ has support contained in $J^\pm(\supp f)$. Accordingly,
 	\begin{equation*}
 		\supp\tilde{E}^\pm_\lambda f\subset
 		\begin{cases}
 			J^\pm(\supp f), & J^\pm(\supp f)\cap K=\varnothing ,\\
 			J^\pm(K\cup \supp f), & \text{otherwise,}
 		\end{cases}
 	\end{equation*}
 	so $\tilde{E}^\pm_\lambda$ are $K$-nonlocal Green operators for $P+ A(\lambda)$. Part $(C)$ is complete.
 	
 	\emph{$5.$ Continuous extension.} Due to~\eqref{eq:inverseidentity}, one has the formula
 	\begin{equation*}
 		\tilde{E}_\lambda^\pm = E^\pm - E^\pm A(\lambda)E^\pm + E^\pm A(\lambda)E^\pm \big(I+ A(\lambda)E^\pm\big)^{-1}A(\lambda)E^\pm,
 	\end{equation*}
 	in which the three terms on the right-hand side have continuous
 	extensions as maps from $H^s_0$ to $H^{s+\beta}_\loc$, $H^{s+\beta+\delta}_\loc$ and $H^{s+\beta+2\delta}_\loc$ respectively. Because $\delta>0$, we deduce that $\tilde{E}^\pm_\lambda$
 	extends continuously to a map $H^s_0\to H_\loc^{s+\beta}$ as required.
 	By the Leibniz rule (Corollary~\ref{cor:Leibniz}), this extension is holomorphic in $\lambda$ with respect to the topology of $\LL_b\big(H^s_0,H_\loc^{s+\beta}\big)$.
 	This proves part (D).
 	
 	
 	\emph{$6.$ Support non-increasing modifications.} Finally, suppose condition $(d')$ holds, so that one has $\supp A(\lambda)f\subset \supp f$ as well as
 	$\supp A(\lambda)f\subset K$ for all $f\in C^\infty$. Then $(d)$ also holds,
 	either as a consequence of Peetre's theorem~\cite{Peetre:1960} or by the
 	following direct argument:
 	if $f\in C^\infty$ vanishes identically on $K$, then the carrier of $A(\lambda)f\in C^\infty_K$ satisfies
 	both
 	\begin{equation*}
 		\carr A(\lambda)f \subset \inte(K)
 	\end{equation*}
 	and
 	\begin{equation*}
 		\carr A(\lambda)f\subset \supp A(\lambda)f\subset \supp f=\overline{\carr f}\subset \overline{M\setminus K} = M\setminus \inte(K)
 	\end{equation*}
 	and we deduce that $\carr A(\lambda)f$ is empty, i.e., $A(\lambda)f$ vanishes identically. Therefore, all the conclusions reached previously still hold.
 	
 	It further follows that
 	\begin{equation*}
 		\supp Y^\pm(\lambda) f =\supp A(\lambda)E^\pm f\subset K\cap \supp E^\pm f\subset K\cap J^\pm(\supp f)
 		\qquad \text{for all} \ f\in C_0^\infty;
 	\end{equation*}
iterating, we have in particular that
 	\begin{equation*}
 		\supp \big(Y^\pm(\lambda)\big)^r f\subset K^\pm_f:= K\cap J^\pm(\supp f) \qquad \text{for all} \ r\in \NN,\ f\in C_0^\infty \ \text{and} \ \lambda\in\CC\setminus S.
 	\end{equation*}	
 	Therefore $\big(Y^\pm(\lambda)\big)^r Y^\pm(\lambda) f\in C^\infty_{K^\pm_f}\subset C^\infty_K$ for all $r\in\NN_0$, $f\in C_0^\infty$, $\lambda\in \CC\setminus S$.
 	Using Corollary~\ref{cor:CKinfty_inversion}, applied to the function $\lambda\mapsto Y^\pm(\lambda) f\in C^\infty_K$, it follows that both
 	$\supp \big(I+Y^\pm(\lambda)\big)^{-1} Y^\pm(\lambda) f$ and
 	\begin{equation*}
 		\supp Y^\pm(\lambda) \big(I+Y^\pm(\lambda)\big)^{-1} Y^\pm(\lambda) f= \supp \big(I-(I+ Y^\pm(\lambda))^{-1}\big) Y^\pm(\lambda) f
 	\end{equation*}
 	are contained in $K^\pm_f$ for all $\lambda\in \CC\setminus S$. Using the identity~\eqref{eq:inverseidentity},
 	it now follows that
 	\begin{equation*}
 		\supp \big(I+ A(\lambda)E^\pm\big)^{-1} f \subset \supp f\cup \left(K\cap J^\pm(\supp f)\right) \subset
 		J^\pm(\supp f)
 	\end{equation*}
 	and hence $\supp \tilde{E}^\pm_\lambda f \subset J^\pm\big(J^\pm(\supp f)\big)=J^\pm(\supp f)$, which is the support property (G3) for standard Green operators, completing the proof of part $(E)$ and therefore the whole theorem.
 \end{proof}


% Original source lines 1284--1534
 \appendix
 \section{Some topological vector spaces}\label{appx:topspaces}

 We briefly rehearse the definition and main properties of the various $C^k$ and Sobolev spaces encountered in the text, broadly following~\cite{Baer:2015,BaerWafo:2015}, before turning to some properties of the topology of bounded convergence that are also needed. No originality is claimed for the material given here.

 \subsection[Spaces of smooth and C\^{}k functions]{Spaces of smooth and $\boldsymbol{C^k}$ functions}
 Let $M$ be a smooth manifold and let $C^k(M)$, $k\in\NN_0\cup\{\infty\}$, be the vector space of complex-valued $k$-times continuously differentiable functions on $M$. For each $k\in\NN_0$ and compact $K\subset M$, one has a seminorm
 \begin{equation*}
 	\|f\|_{K,k} = \max_{0\le r\le k} \max_{x\in K} |(\nabla^r f)(x)|_r
 \end{equation*}
 on $C^k(M)$, where $\nabla$ is an arbitrarily chosen connection on $M$ and $|\cdot|_r$ an arbitrarily chosen norm making $T^*M^{\otimes r}$ a (finite-dimensional) Banach bundle; different choices result in equivalent seminorms. The collection of seminorms $\|\cdot\|_{K,k}$ as $K$ runs over compact subsets of $M$ and $k\in\NN_0$ provides a Fr\'echet topology on $C^\infty(M)$; similarly,
 we obtain a Fr\'echet topology on~$C^k(M)$, $k\in\NN_0$, using the seminorms
 $\|\cdot\|_{K,k}$.


 If $A$ is closed, we define $C^\infty_A(M)=\{f\in C^\infty(M)\colon\supp f\subset A\}$ with the relative topology. Thus the topology is defined by the seminorms $\|\cdot\|_{K,k}$ as $K$ runs over compact subsets $K\subset A$ and $k\in\NN_0$; if $A$ is compact, it is sufficient to use the seminorms $\|\cdot\|_{A,k}$, $k\in\NN_0$. As $C^\infty_A(M)$ is closed subspace of a Fr\'echet space, it is also Fr\'echet. Defining $C^k_A(M)$ as the analogous subspace of $C^k(M)$, the topology is generated by seminorms $\|\cdot\|_{K,k}$ for compact $K\subset A$, or just the single seminorm $\|\cdot\|_{A,k}$ (which is a norm on $C^k_A(M)$) in the case that $A$ is compact.

 A \emph{support system}~\cite{Baer:2015} is a subset $\mathscr A$ of the set of all closed subsets on $M$, which is closed under finite unions and has the property that for each $A\in \mathscr A$ it holds that $(i)$
 $A\subset \textrm{int}\,(A')$ for some $A'\in\mathscr A$ and $(ii)$ if $A'$ is a closed subset of $M$ with $A'\subset A$ then $A'\in\mathscr A$. Any support system is a directed system with respect to inclusion and we write
 \begin{equation*}
 	C^\infty_{\mathscr{A}}(M)= \bigcup_{A\in\mathscr{A}} C^\infty_A(M)
 \end{equation*}
 with the locally convex inductive limit topology,
 so that a convex set $U\subset C^\infty_{\mathscr{A}}(M)$ is a neighbourhood of $0$
 if and only if $U\cap C^\infty_A(M)$ is a neighbourhood of $0$ in $C^\infty_A(M)$ for every $A\in\mathscr{A}$; because $\mathscr{A}$ is directed, one also has that
 a convex set $O$ is open if and only if $O\cap A$ is open in each $C^\infty_A(M)$.

 Examples of support systems include the set of compact sets, leading to the space of compactly supported functions $\CoinX{M}$, and the sets of (strictly) future/past/spatially-compact sets, giving rise to $C^\infty_{\sfc/\spc/\fc/\pc/\mathrm{sc}}(M)$.

 A linear map $T$ from $C^\infty_{\mathscr{A}}(M)$ to a locally convex topological vector space $Y$ is continuous if and only if all its restrictions $T_A\colon C^\infty_A(M)\to Y$ are continuous ($A\in\mathscr{A}$). In particular, for a linear map $T\colon C^\infty_{\mathscr{A}}(M)\to C^\infty_{\mathscr{B}}(M)$ to be continuous, it suffices that each restriction $T_A$ has range contained in some $C^\infty_B(M)$ for some $B\in\mathscr{B}$ (depending on $A$) and determines a continuous map $C^\infty_A(M)\to C^\infty_B(M)$, thus implying that $T_A\colon C^\infty_A(M)\to C^\infty_{\mathscr{B}}(M)$ is continuous. We record the following application:
 \begin{Lemma}%\label{lem:restr}
 	If $T$ is a continuous linear self-map of $C^\infty(M)$ such that $\supp Tf\subset K\cup \supp f$ for all $f\in C^\infty(M)$, where $K$ is a fixed compact set, then $T$ restricts to any of $C^\infty_{0/\sfc/\spc/\fc/\pc/\mathrm{sc}}(M)$ as a continuous map.
 \end{Lemma}
 \begin{proof}
 	If $\mathscr{A}$ is one of the relevant support systems then $A\cup K\in\mathscr{A}$ for each $A\in\mathscr{A}$. As the restriction $T_A$ of $T$ to $C^\infty_A(M)$ has its range in $C^\infty_{A\cup K}(M)$, and both these function spaces are subspaces of $C^\infty(M)$ with the relative topology, continuity of $T_A\colon C^\infty_A(M)\to C^\infty_{A\cup K}(M)$ follows from that of $T$. Letting $A$ vary in $\mathscr{A}$, the result is proved.
 \end{proof}

 Letting $\Omega_\alpha$ be the bundle of densities of weight $\alpha$, we can define spaces $\Gamma^\infty(\Omega_\alpha)$,
 $\Gamma_A^\infty(\Omega_\alpha)$ and $\GoinX{\Omega_\alpha}$ of smooth sections
 of $\Omega_\alpha$ in a similar way; any smooth nowhere vanishing density $\mu$ on
 $M$ provides topological isomorphisms between these spaces and their zero-weight analogues, simply by multiplication by appropriate powers of $\mu$. The space of distributions on $M$, $\DD'(M)$ is the topological dual of $\GoinX{\Omega_1}$,
 equipped with the weak-$*$ topology. In particular, $C^\infty(M)$ is canonically embedded as a subspace of $\DD'(M)$. More generally,
 $\DD'(M,\Omega_\alpha)$ is defined as the dual of $\GoinX{\Omega_{1-\alpha}}$.
 We also write $\EE'(M)$ for the topological dual of $\Gamma^\infty(\Omega_1)$, i.e.,
 the distributions of compact support.


 \subsection{Sobolev spaces}

 \subsubsection{Compact manifolds}

 {\bf The spaces $\mathbf{H^s(M)}$.}
 On a \emph{compact} manifold $M$, choose an auxiliary Riemannian metric and define $L^2(M)$ in the usual way, using the volume element induced by the metric and writing the inner product as $\langle\cdot\mid\cdot\rangle$ (linear in the second slot). Let $T=(-\triangle + I)^{1/2}$, where $\triangle$ is the Laplace--Beltrami operator, initially defined on $C^\infty(M)$ and then extended (uniquely) to a self-adjoint negative operator on $L^2(M)$.
 Its complex powers $T^s$ (as defined by functional calculus) are classical pseudodifferential operators of order $s$, $T^s\in\Psi^s(M)$~\cite{Seeley:1967} (see~\cite{Amman_etal:2004} for an axiomatically based proof) and $T^s$ is compact for $\Re s<0$ (most easily seen using the spectral properties of~$-\triangle$ on compact manifolds). The domain of $T^s$ contains $C^\infty(M)$ for all $s$.

 For $s\in\RR$, the Sobolev space $H^s(M)$ is defined as the completion of $C^\infty(M)$ with respect to the norm
 \begin{equation*}
 	\|f\|_{H^s(M)} = \| T^s f\|_{L^2(M)}
 \end{equation*}
 and, as a topological vector space, is independent of the choice of auxiliary Riemannian metric involved in its construction (of course, the specific norm $\|\cdot\|_{H^s(M)}$
 and its compatible Hilbert space inner product
 are metric-dependent). Indeed, any positive second-order elliptic operator that is essentially self-adjoint on $C^\infty(M)$ could be used in place of $-\triangle$.
 Evidently $T^s$ extends to an isometry from $H^s(M)$ to $L^2(M)$, which may be used to embed $H^s(M)$ in $\DD'(M)$ so that $u \in H^s(M)$ corresponds to the distribution $f\mapsto \ip{T^{-s}\overline{f/\nu}}{T^su}$, $f\in\GoinX{\Omega_1}$, where $\nu$ is the density corresponding to the
 volume element of the auxiliary metric used to define $L^2(M)$. Restricted to $u\in C^\infty(M)$, this embedding is consistent with the embedding of $C^\infty(M)$ in $\DD'(M)$ already mentioned.
 As $T^q$ is compact on $L^2(M)$ for $q<0$, it follows that $H^s(M)\subset H^t(M)$ for all $s>t$, with a compact inclusion map. More generally. $P\colon H^{s+m}(M)\to H^s(M)$ is continuous for any partial differential operator of order $m$ with smooth coefficients (or indeed any pseudodifferential operator $P\in \Psi^{m}(M)$), because $T^{-m}P\in \Psi^0(M)$ extends to a bounded operator on $L^2(M)$.


 {\bf Duality.} For any $s\in\RR$, suppose that $\ell\in H^s(M)'$, so there is a constant $c$ so that $|\ell(u)|\le c\|u\|_{H^s(M)}$ for all $u\in H^s(M)$. Then $|\ell(T^{-s}f)|\le c\|f\|_{L^2(M)}$ for $f\in L^2(M)$ and consequently there is $w\in L^2(M)$ so that $\ell(u) =\langle w\,|\, T^{s}u\rangle_{L^2(M)}$ for all $u\in H^s(M)$. Choose a sequence $w_n\to w$ in $L^2(M)$ with
 $w_n\in C^\infty(M)$ and note that $T^s w_n\in C^\infty(M)$ is a Cauchy sequence with respect to the
 $H^{-s}(M)$ norm, converging to some $v\in H^{-s}(M)$ for which
 $T^{-s}v=\lim_n w_n=w$. Thus
 $\ell(u)=\langle T^{-s}v\, |\,T^su\rangle$ for all $u\in H^s(M)$ and
 $|\ell(u)|\le \|v\|_{H^{-s}(M)}\|u\|_{H^s(M)}$. Noting that
 \begin{equation*}
 	\frac{\ell(T^{-s}w_n)}{\|T^{-s}w_n\|_{H^s(M)}}=\frac{\langle w\mid w_n \rangle_{L^2(M)}}{\|w_n\|_{L^2(M)}}\to \|w\|_{L^2(M)} = \|v\|_{H^{-s}(M)},
 \end{equation*}
 we see that $\|\ell\|=\|v\|_{H^{-s}(M)}$, and we have established that
 the space $H^{-s}(M)$ is anti-isomorphic to $H^s(M)'$ with the operator norm topology (i.e.,
 the strong dual) with respect to the sesquilinear pairing $\langle v,u\rangle = \langle T^{-s}v\,|\, T^s u\rangle$, $v\in H^{-s}(M)$, $u\in H^s(M)$. To be precise: every
 choice of Riemannian metric on $M$ induces an anti-isomorphism of this type, and the
 pairing depends on the choice made.

 {\bf Embedding theorems.}
 The relationship between $C^k$ and Sobolev spaces is given as follows. In one direction, the formula $-\triangle=\delta d$ (on $0$-forms) gives the estimates
 \begin{equation*}
 	0\le \ip{f}{(-\triangle)^{2r}f} = \| (-\triangle)^r f\|^2_{L^2(M)}\le C\|f\|_{M,2r}^2
 \end{equation*}
 and
 \begin{equation*}
 	0\le \ip{f}{(-\triangle)^{2r+1}f} = \| d(-\triangle)^r f\|^2_{\Lambda^1(M)}\le C\|f\|_{M,2r+1}^2 ,
 \end{equation*}
 where $\Lambda^1(M)$ is the Hilbert space of square-integrable $1$-forms (with respect to the auxiliary Riemannian metric) and $f\in C^\infty(M)$.
 From this we find
 \begin{equation*}
 	\|f\|_{H^k(M)}^2 = \ip{f}{(-\triangle+I)^k f} = \sum_{j=0}^k \binom{k}{j} \ip{f}{(-\triangle)^{j}f} \le C \|f\|_{M,k}^2
 \end{equation*}
 for all $k\in \NN_0$. Thus $C^k(M)$ is continuously embedded in $H^s(M)$ for all $s\ge k$.

 On the other hand, now let $n$ be the maximum dimension of any component of $M$.
 For $\Re s>n/2$, the integral kernel $Z_s(p,q)$ of the operator $(-\triangle+I)^{-s}=T^{-2s}$ on $L^2(M)$ is continuous, $Z_s\in C(M\times M)$ and can also be written
 in terms of the spectral decomposition of
 $(-\triangle+I)$ as
 \begin{equation*}
 	Z_s(p,q) = \sum_{r} \frac{e_r(p)\overline{e_r(q)}}{\lambda_r^{s}},
 \end{equation*}
 where $e_r\in L^2(M)$ are a basis of smooth orthonormal eigenfunctions for $-\triangle+I$ with corresponding eigenvalues $\lambda_r$ -- see~\cite{Seeley:1967}. In particular,
 \begin{equation*}
 	v_p = \sum_{r} \frac{\overline{e_r(p)}e_r}{\lambda_r^{s/2}}
 \end{equation*}
 belongs to $L^2(M)$ with $\|v_p\|^2 = Z_s(p,p)$ by the Pythagoras theorem.
 For $\Re s>n/2$, it follows that $\|e_r\|_\infty \le C_s\big|\lambda_r^{s/2}\big|$, where $C_s=\sup_{p\in M}|Z_s(p,p)|^{1/2}$. Moreover, if $f\in C^\infty(M)$,
 \begin{equation*}
 	|\langle e_r \mid f \rangle|=\lambda_r^{-k}\big|\ip{e_r}{T^k f}\big|\le \lambda^{-k}_r\|f\|_{H^k(M)},\qquad k\in\NN,
 \end{equation*}
 so $\ip{e_r}{f}$ decays faster than any inverse power of $\lambda_r$ and
 $\ip{v_p}{T^s f} = \sum_r e_r(p)\ip{e_r}{f}= f(p)$ for all (not merely almost all) $p\in M$. Consequently,
 \begin{equation*}
 	\|f\|_\infty =\sup_{p\in M} |\langle v_p \mid f \rangle|\le C_s \|f\|_{H^s(M)}, \qquad f\in C^\infty(M).
 \end{equation*}
 By density of $C^\infty(M)$ in $H^s(M)$, it follows
 that there is a continuous embedding $H^s(M)\to C(M)$ if $s>n/2$; considering $\|Pf\|_\infty$ in a~similar way for differential operators $P$, one sees that $H^s(M)$ is continuously embedded in $C^k(M)$ for all $s>k+n/2$. It follows that $\cap_{s\in\RR} H^s(M)=C^\infty(M)$.


 If $u\in \DD'(M)= \GoinX{\Omega_1}'$ is a distribution then, owing to compactness of $M$, there is a~partial differential operator $P$ so that $|u(\nu f)|\le \|Pf\|_\infty \le C_s \|Pf\|_{H^s(M)}$ for any $s>n/2$ and all $f\in C^\infty(M)$. Thus, with $t=s+\operatorname{ord} (P)$, $u(\nu f)= \big\langle \overline{f},w\big\rangle= \ip{T^t \overline{f}}{T^{-t}w}_{L^2(M)}$ for some $w\in H^{-t}(M)$, which is evidently compatible with the embedding of $H^{-t}(M)$ in $\DD'(M)$
 described earlier. In this sense,
 \begin{equation*}
 	\DD'(M) = \bigcup_{s\in\RR} H^{s}(M).
 \end{equation*}

 \subsubsection{General manifolds}

 For a general (not-necessarily compact) smooth manifold $M$ with at most finitely many components,
 we proceed as follows. If $K\subset M$ is compact and topologically regular,\footnote{Note that if $K$ is any compact subset, then the closure of the interior of $K$ is compact and topologically regular.} we may
 diffeomorphically identify $K$ with a compact subset $\hat{K}$ of a compact manifold $N$ (for example, by `doubling' a compact set that contains $K$ in its interior and has a smooth boundary~\cite{BaerWafo:2015}; for~$K$ contained in a coordinate chart one could take $N$ to be a torus). Then the Sobolev space~$H^s_K(M)$ is defined as the completion of
 $C^\infty_K(M)$ with respect to the pull back of (some choice of) the~$H^s(N)$ norm. The space $H^s_K(M)$ can be identified with a subspace of $\DD'_K(M)$, the distributions with support contained in $K$. As a topological vector space, it is independent of the choices used in its construction. However, by making such a choice one may
 endow $H^s_K(M)$ with a norm, denoted $\|\cdot\|_{H^s_K(M)}$, and indeed a compatible Hilbert space structure.
 Estimates of the form
 $\|f\|_{H^k_K(M)}\le C\|f\|_{K,k}$, $f\in C^\infty_K(M)$, and hence continuity of
 the embedding $C^k_K\to H^k_K$, carry over from the compact case for each fixed $k\in\NN_0$, as does compactness of the embedding $H^{s}_K(M)\to H^{t}_K(M)$ for $s>t$
 and the continuous embedding of $H^s_K(M)$ in $C^k_K(M)$ for $s>k+n/2$,
 where again $n$ is the maximum dimension of any component of $M$.
 Consequently one has $\cap_{s\in\RR} H^s_K(M)=C^\infty_K(M)$.

 Next, we define (abbreviating `compact and topologically regular' by c.t.r.)
 \begin{equation*}
 	H_0^s(M) = \bigcup_{\text{c.t.r. $K\subset M$}} H^s_K(M),
 \end{equation*}
 with the locally convex inductive limit topology.
 As $M$ admits countable compact exhaustions (and consequently, countable c.t.r.\ exhaustions), $H_0^s(M)$ may be realised as a countable strict inductive limit -- i.e., it is a $LF$ space. A fact to be used later is that, by~\cite[Proposition~14.6]{Treves:TVS} (see also Proposition~4 in~\cite{DieudonneSchwartz:1949}) a subset $B$ of $H_0^s(M)$ is bounded if and only if $B$ is a bounded subset of~$H^s_K(M)$ for some compact $K$. It is easily shown that
 \begin{equation*}
 	\EE'(M) = \bigcup_{s\in\RR} H_0^{s}(M).
 \end{equation*}






 Finally, the local Sobolev space $H^s_\loc(M)$ is defined as
 \begin{equation*}
 	H^s_\loc(M) = \big\{u\in\DD'(M)\colon \text{$\chi u\in H^s_0(M)$ for all $\chi\in C_0^\infty({M})$}\big\},
 \end{equation*}
 and is equipped with the Fr\'echet topology induced by the seminorms $\|\chi \cdot\|_{H^s_{K}(M)}$ as $K$ runs over compact topologically regular subsets of $M$ and $\chi$ runs over $C^\infty_K(M)$. The inclusion
 $H^s_0(M)\hookrightarrow H^s_\loc(M)$ is continuous for all $s$ and indeed we have
 \begin{equation*}
 	H^s_0(M) = H^s_\loc(M) \cap \EE'(M).
 \end{equation*}
 Thus, if $M$ is compact, $H^s_\loc(M)$ and $H^s_0(M)$ coincide as topological vector spaces. The construction above produces the same spaces as the chart-based approach taken in~\cite{Hormander3}.

 The Sobolev embedding theorems already mentioned entail the existence of continuous embeddings of $C^k(M)$ in $H^k_\loc(M)$ for all $k\in\NN_0$ and of $H^s_\loc(M)$ in $C^k(M)$ for all $s>k+n/2$, where~$n$ is the maximum dimension of any component of $M$. Consequently, $\cap_{s\in\RR} H^s_\loc(M)=C^\infty(M)$.

 \subsection{The topology of bounded convergence}\label{appx:top_bded_cvgnce}

 A general reference for the following is Tr\`eves~\cite[Chapter~32]{Treves:TVS}.

 If $E$ and $F$ are Hausdorff locally convex topological spaces then the topology of bounded convergence on $\LL(E,F)$, the space of all continuous linear maps $E\to F$, is defined by the neighbourhood base of zero, consisting of sets
 \begin{equation*}
 	\Uf(B;V) = \{T\in \LL(E,F)\colon T(B)\subset V\}
 \end{equation*}
 as $B$ runs over the bounded subsets of $E$ and $V$ runs over any neighbourhood base of zero in the topology of $F$. The notation $\LL_b(E,F)$ denotes $\LL(E,F)$ equipped with the topology of bounded convergence.
 Thus a net $T_\alpha$ converges to $0$ in $\LL_b(E,F)$
 if and only if, for every bounded $B\subset E$ and neighbourhood base set $V$,
 $T_\alpha$ eventually maps $B$ into $V$. This topology makes $\LL(E,F)$ Hausdorff and locally convex (inherited from $F$~\cite[p.~336]{Treves:TVS}). We record some basic facts.
 \begin{Lemma}\label{lem:Lb_basic}
 	Let $E$, $F$, $G$ be Hausdorff locally convex topological spaces. $(a)$ If $T\in \LL(E,F)$ and the net $S_\alpha\to S$ in $\LL_b(F,G)$ then $S_\alpha T\to ST$ in $\LL_b(E,G)$; $(b)$ if the net $T_\alpha\to T$ in $\LL_b(E,F)$ and $S\in \LL(F,G)$ then
 	$ST_\alpha\to ST$ in $\LL_b(E,G)$; $(c)$ if $T_\alpha\to T$ in $\LL_b(E,F)$ and $\Ran T_\alpha\subset \hat{F}$, where $\hat{F}$ is a topological subspace of $F$, then $T_\alpha\to T$ in $\LL_b(E,\hat{F})$; $(d)$ if $F$ is barrelled and $T_n\to T$ and $S_n\to S$ are convergent \emph{sequences} in $\LL_b(E,F)$ and $\LL_b(F,G)$, then $S_n T_n\to ST$ in~$\LL_b(E,G)$.
 \end{Lemma}
 Note that Hilbert, Banach, Fr\'echet and LF spaces are all barrelled~\cite[Chapter~33]{Treves:TVS}.
 \begin{proof}
 	$(a)$ It is enough to prove this in the case $S=0$; taking any bounded $B$ in $E$ with $0\in B$, and any $0$-neighbourhood $V$ in $G$, note that $T(B)$ is bounded in $F$ and deduce that $S_\alpha T(B)$ is eventually contained in $V$. Thus $S_\alpha T\to 0$ in $\LL_b(E,G)$. $(b)$ Without loss, suppose $T=0$ and with $B$ and $V$ as before, we note that $S^{-1}(V)$ is a $0$-neighbourhood in $F$ and again deduce that $ST_\alpha(B)$ is eventually contained in $V$, so $ST_\alpha\to 0$.
 	
 	$(c)$ Again, it is enough to treat $T=0$. Take any bounded $B\subset E$ and $0$-neighbourhood~$\hat{V}$ in~$\hat{F}$. Then $\hat{V}=\hat{F}\cap V$ for an $0$-neighbourhood~$V$ in $F$ because $\hat{F}$ carries the subspace topology. We know that, eventually, $T_\alpha(B)\subset V$, and as $\Ran T_\alpha\subset \hat{F}$, we have, eventually, that $T_\alpha(B)\subset V\cap\hat{F}= \hat{V}$. Thus $T_\alpha\to 0$ in $\LL_b(E,\hat{F})$.
 	
 	$(d)$ Suppose first that $T=0$ and let $0\in B$
 	be any bounded subset of $E$ and $V$ any $0$-neighbourhood in $G$. Then
 	the $S_n$, together with $S$, form a bounded subset in $\LL_b(F,G)$ (which, we recall, is Hausdorff). As $F$ is barrelled, they form an equicontinuous set of maps by the Banach--Steinhaus theorem~\cite[Theorem 33.1]{Treves:TVS}. Thus there exists
 	a $0$-neighbourhood $W$ in $F$ such that $S_n (W)\subset V$ for all $n$.
 	As $T_n(B)\subset W$ for all sufficiently large $n$, we have $S_nT_n(B)\subset W$ for such $n$. Thus $S_nT_n\to 0$ in this case. If $T\neq 0$, we now know
 	that $S_n(T_n-T)\to 0$, and as $S_n T\to ST$ we find $S_nT_n\to ST$ as required.
 \end{proof}

 Typical applications of Lemma~\ref{lem:Lb_basic}\,$(a)$,\,$(b)$ are to show, for instance, that a holomorphic function from a domain in $\CC$ to $\LL_b\big(H^s_0(M),H^s_K(M)\big)$ is also holomorphic as a function to $\LL_b\big(H^s_K\big)$ and
 $\LL_b(H^s_0(M))$ due to continuous embedding of $H_K^s(M)$ in $H^s_0(M)$. Part~$(d)$ has the following use.\looseness=-1
 \begin{Corollary}\label{cor:Leibniz}
 	Let $E$, $F$, $G$ be Hausdorff locally convex topological spaces, with $F$ being barrelled.
 	If $\lambda\mapsto T(\lambda)$ and $\lambda\mapsto S(\lambda)$ are functions from a domain in $\CC$ to~$\LL(E,F)$ and~$\LL(F,G)$ respectively, then $(a)$ if $T$ and $S$ are both continuous in the topology of bounded convergence $($at $\mu$$)$, so is $(ST)(\lambda)=S(\lambda)T(\lambda)$ $($at $\mu$$)$; $(b)$ if $T$ and $S$ are differentiable at $\mu$ in the topology of bounded convergence then so is $ST$, with a derivative given by the Leibniz rule
 	$(ST)'(\mu)=S'(\mu)T(\mu) + S(\mu)T'(\mu)$.
 \end{Corollary}
 \begin{proof}
 	As $\CC$ is first-countable, questions of continuity and differentiability may be reduced to sequential considerations, and the result follows using Lemma~\ref{lem:Lb_basic}\,$(d)$ and the standard proof of the Leibniz formula.
 \end{proof}

 It is also useful to have some sufficient conditions for convergence in $\LL_b(E,F)$ topology where~$E$ and $F$ are Fr\'echet or countable strict inductive limits thereof (LF spaces). If $F$ is Fr\'echet, then the bounded subsets are precisely the subsets $B$ so that $\sup_{f\in B} \rho_j(f)<\infty$ for every seminorm~$\rho_j$ defining the topology of $F$, while sets $V_{j,\epsilon}=\{f\in F\colon\rho_j(f)<\epsilon\}$, for arbitrary~$\epsilon>0$ and defining seminorm $\rho_j$, provide a basis of neighbourhoods of zero. On the other hand, if $E$ is an LF space with defining sequence $E_n$ of Fr\'echet spaces, then
 the bounded subsets~$B$ of~$E$ comprise precisely those subsets that are bounded subsets of some $E_n$~\cite[Proposition~14.6]{Treves:TVS}, while a neighbourhood base is provided by convex sets $V\subset E$ so that
 each~$V\cap E_n$ is an open neighbourhood of zero in $E_n$.
 We recall that continuous linear maps between topological vector spaces preserve boundedness~\cite[Proposition~14.2]{Treves:TVS}. The following results are used in the text.
 \begin{Lemma} \label{lem:LbFrechet}
 	Suppose $T_\alpha\in \LL_b(F)$ is a net of operators on a Fr\'echet space $F$ with defining seminorms $\rho_j$. If, for each $j$, there exists $k(j)$,
 	such that for all $\epsilon>0$, it is eventually true that $\rho_j(T_\alpha f)<\epsilon\rho_{k(j)}(f)$ for all $f\in F$, then $T_\alpha\to 0$ in $\LL_b(F)$.
 \end{Lemma}
 \begin{proof}
 	Given any bounded set $B$ and
 	neighbourhood $V_{j,\epsilon}$, we set $C=\sup_{f\in B}\rho_{k(j)}(f)$ and apply the given property to $\epsilon/C$ to find that, eventually,{\samepage
 	\begin{equation*}
 		\rho_j(T_\alpha f)<\epsilon C^{-1} \rho_{k(j)}(f) \le \epsilon, \qquad \forall f\in B.
 	\end{equation*}
 	We have shown that, eventually, $T_\alpha\in \Uf(B;V_{j,\epsilon})$; as $j$ and $\epsilon$ were arbitrary, $T_\alpha\to 0$ in $\LL_b(F)$.}
 \end{proof}

 \begin{Corollary} \label{cor:HstoCinfty}
 	Consider a net $T_\alpha\in \LL\big(C^\infty_K\big)$
 	such that each $T_\alpha$ extends to an operator in~$\LL(H^s_K)$ $($which we also denote $T_\alpha$$)$ for all $s\ge s_*$. Suppose also that $T_\alpha\to 0$ in every $\LL(H^s_K)$, $s\ge s_*$. Then $T_\alpha\to 0$ in $\LL_b\big(C^\infty_K\big)$.
 \end{Corollary}
 \begin{proof}
Any $C^\infty_K$ seminorm is bounded by an $H^s_K$ norm and vice versa, so, for any
 	$C^\infty_K$-seminorm $\|\cdot\|_{K,j}$ we may estimate
 	\begin{equation*}
 		\|T_\alpha f\|_{K,j} \le c \| T_\alpha f\|_{H^{s(j)}_K} \le c\|T_\alpha\|_{\LL(H^{s(j)}_K)} \|f\|_{H^{s(j)}_K}
 		\le c' \|T_\alpha\|_{\LL(H^{s(j)}_K)} \|f\|_{K,k(s(j))}
 	\end{equation*}
 	for suitable choices of $s(j)$ and $k(s(j))$, uniformly in $f\in C^\infty_K$ and $\alpha$.
 	As $T_\alpha\to 0$ in $\LL\big(H^{s(j)}_K\big)$, for any $\epsilon>0$ it is eventually true that
 	$\|T_\alpha f\|_{K,j} \le \epsilon \|f\|_{K,k(s(j))}$ for all $f\in C^\infty_K$. Hence $T_\alpha\to 0$ in $\LL_b\big(C^\infty_K\big)$ by Lemma~\ref{lem:LbFrechet}.
 \end{proof}

 \begin{Lemma}\label{lem:LbLF}
 	Suppose $T_\alpha$ is a net of operators on the LF space $E=\bigcup_{n\in\NN} E_n$. If for some fixed $n$, one has $\Ran(T_\alpha)\subset E_n$ and $T_\alpha|_{E_m}\to 0$ in $\LL_b(E_m,E_n)$ for all $m$, then $T_\alpha\to 0$ in both $\LL_b(E,E_n)$ and $\LL_b(E)$.
 \end{Lemma}
 \begin{proof}
 	Consider any neighbourhood $V$ of zero in the standard neighbourhood base of $E_n$ and any bounded set $B\subset E$, necessarily obeying $B\subset E_m$ for some $m$. Then $T_\alpha$ eventually maps $B$ into $V$; as $B$ and $V$ were arbitrary, $T_\alpha\to 0$ in $\LL_b(E,E_n)$. Post-composing with the continuous embedding of $E_n$ in $E$, $T_\alpha\to 0$ in $\LL_b(E)$ as well.
 \end{proof}
\begin{thebibliography}{99}
\bibitem[1]{Amman_etal:2004}
Ammann B., Lauter R., Nistor V., Vasy A., Complex powers and non-compact
 manifolds, \href{https://doi.org/10.1081/PDE-120037329}{\textit{Comm. Partial Differential Equations}} \textbf{29} (2004),
 671--705, \href{https://arxiv.org/abs/math.OA/0211305}{arXiv:math.OA/0211305}.

\bibitem[2]{Baer:2015}
B\"ar C., Green-hyperbolic operators on globally hyperbolic spacetimes,
 \href{https://doi.org/10.1007/s00220-014-2097-7}{\textit{Comm. Math. Phys.}} \textbf{333} (2015), 1585--1615,
 \href{https://arxiv.org/abs/1310.0738}{arXiv:1310.0738}.

\bibitem[4]{BaerWafo:2015}
B\"ar C., Wafo R.T., Initial value problems for wave equations on manifolds,
 \href{https://doi.org/10.1007/s11040-015-9176-7}{\textit{Math. Phys. Anal. Geom.}} \textbf{18} (2015), 7, 29~pages,
 \href{https://arxiv.org/abs/1408.4995}{arXiv:1408.4995}.

\bibitem[5]{Bernal:2004gm}
Bernal A.N., S\'{a}nchez M., Smoothness of time functions and the metric
 splitting of globally hyperbolic spacetimes, \href{https://doi.org/10.1007/s00220-005-1346-1}{\textit{Comm. Math. Phys.}}
 \textbf{257} (2005), 43--50, \href{https://arxiv.org/abs/gr-qc/0401112}{arXiv:gr-qc/0401112}.

\bibitem[6]{Bernal:2006xf}
Bernal A.N., S\'{a}nchez M., Globally hyperbolic spacetimes can be defined as
 `causal' instead of `strongly causal', \href{https://doi.org/10.1088/0264-9381/24/3/N01}{\textit{Classical Quantum Gravity}}
 \textbf{24} (2007), 745--749, \href{https://arxiv.org/abs/gr-qc/0611138}{arXiv:gr-qc/0611138}.

\bibitem[9]{DieudonneSchwartz:1949}
Dieudonn\'e J., Schwartz L., La dualit\'e dans les espaces {$\mathcal{F}$} et
 {$(\mathcal{LF})$}, \href{https://doi.org/10.5802/aif.8}{\textit{Ann. Inst. Fourier (Grenoble)}} \textbf{1} (1949),
 61--101.

\bibitem[14]{FewVer:20xx}
Fewster C.J., Verch R., A nonlocal generalization of {G}reen hyperbolicity,
 {i}n preparation.

\bibitem[15]{Geroch_spinorI:1968}
Geroch R., Spinor structure of space-times in general relativity.~{I},
 \href{https://doi.org/10.1063/1.1664507}{\textit{J.~Math. Phys.}} \textbf{9} (1968), 1739--1744.

\bibitem[18]{Hormander1}
H\"ormander L., The analysis of linear partial differential operators.~{I}.
 {D}istribution theory and Fourier analysis, \textit{Classics Math.}, \href{https://doi.org/10.1007/978-3-642-61497-2}{Springer}, Berlin,
 2003.

\bibitem[19]{Hormander3}
H\"ormander L., The analysis of linear partial differential operators.~{III}.
 Pseudo-differential operators, \textit{Grundlehren Math. Wiss.}, Vol.~274,
 \href{https://doi.org/10.1007/978-3-540-49938-1}{Springer}, Berlin, 2007.

\bibitem[25]{Peetre:1960}
Peetre J., R\'ectification \`a l'article ``{U}ne caract\'erisation abstraite
 des op\'erateurs diff\'erentiels'', \href{https://doi.org/10.7146/math.scand.a-10598}{\textit{Math. Scand.}} \textbf{8} (1960),
 116--120.

\bibitem[26]{ReedSimon:vol1}
Reed M., Simon B., Methods of modern mathematical physics. {I}.~{F}unctional
 analysis, 2nd ed., Academic Press, New York, 1980.

\bibitem[27]{Seeley:1967}
Seeley R.T., Complex powers of an elliptic operator, in Singular {I}ntegrals
 ({P}roc. {S}ympos. {P}ure {M}ath., {C}hicago, {I}ll., 1966), \href{https://doi.org/10.1090/pspum/010/0237943}{American Mathematical Society}, Providence, R.I., 1967, 288--307.

\bibitem[28]{Treves:TVS}
Tr\`eves F., Topological vector spaces, distributions and kernels, Dover
 Publications, Inc., Mineola, NY, 2006.

\end{thebibliography}
\end{document}
